\subsection{Cycle 1-codimensionnel associé à un diviseur}
\label{subsection:IV.21.6}
\label{IV.21.6}

\begin{env}[21.6.1]
\label{IV.21.6.1}
Soit $X$ un préschéma localement noethérien, et soit $\mathfrak{I}(X)$ l'ensemble des \emph{parties fermées irréductibles} de $X$ (qui est en correspondance biunivoque avec $X$ par l'application $x\mapsto\overline{\{x\}}$).
Dans le groupe produit $\bb{Z}^X$, considérons le sous-groupe $\mathfrak{Z}(X)$ des éléments $(n_x)_{x\in X}$ tels que l'ensemble des $\overline{\{x\}}\in\mathfrak{I}(X)$ tels que $n_x\neq0$ (ou, ce qui revient au même, l'ensemble des $x\in X$ tels que $n_x\neq0$) soit localement fini.
Il est clair que $\mathfrak{Z}(X)$ est un sous-groupe de $\bb{Z}^X$ qui contient le groupe somme directe $\bb{Z}^{(X)}$ (groupe libre ayant pour base $\mathfrak{I}(X)$), et lui est égal lorsque $X$ est noethérien.
Les éléments de $\mathfrak{Z}(X)$ sont appelés \emph{cycles} sur $X$ et ceux de $\mathfrak{I}(X)$ \emph{cycles premiers} (ils ne forment pas en général une base de $\mathfrak{Z}(X)$ lorsque $X$ n'est pas noethérien).
On considère toujours $\mathfrak{Z}(X)$ comme ordonné par l'ordre induit sur ce sous-groupe par l'ordre produit de $\bb{Z}^X$ et on notera $\mathfrak{Z}^+(X)$ l'ensemble des cycles $\geq0$.

\oldpage[IV-4]{271}

Pour tout cycle $Z\in\mathfrak{Z}(X)$, égal à $(n_x)_{x\in X}$, on écrira par abus de langage,
\[
  Z=\sum_{x\in X}n_x\cdot\overline{\{x\}}.
\]
On appelle \emph{multiplicité de $Z$ au point $x$} et on note $\operatorname{mult}_x(Z)$ l'entier $n_x$ (positif ou négatif).
Dire que $Z\geq0$ signifie que $\operatorname{mult}_x(Z)\geq0$ pour tout $x\in X$.
On appelle \emph{support de $Z$} et on note $\Supp(Z)$ la réunion des $\overline{\{x\}}$ tels que $\operatorname{mult}_x(Z)\neq0$;
par définition de $\mathfrak{Z}(X)$, c'est une partie fermée de $X$, comme réunion d'une famille localement finie de parties fermées.
On appelle \emph{dimension} (resp.\ \emph{codimension}) de $Z$ et on note $\dim(Z)$ (resp.\ $\codim(Z)$) la dimension (resp.\ codimension dans $X$) de $\Supp(Z)$.
\end{env}

\begin{env}[21.6.2]
\label{IV.21.6.2}
On dit qu'une partie fermée $Y$ de $X$ est \emph{purement de codimension $p$} (dans $X$) si \emph{chacune} de ses composantes irréductibles est de codimension $p$ dans $X$.
On dit qu'un cycle est \emph{purement de codimension $p$}, ou (par abus de langage) est un \emph{cycle $p$-codimensionnel}, si son support est purement de codimension $p$.
On note $X^{(p)}$ l'ensemble des $x\in X$ tels que $\codim(\overline{\{x\}},X)=p$, ou, ce qui revient au même \sref{IV.5.1.2}, $\dim(\sh{O}_{X,x})=p$:
les cycles purement de codimension $p$ forment un sous-groupe $\mathfrak{Z}^p(X)$ de $\mathfrak{Z}(X)$, isomorphe au sous-groupe de $\bb{Z}^{X^{(p)}}$ formé des $(n_x)_{x\in X^{(p)}}$ tels que l'ensemble des $\overline{\{x\}}$ (ou l'ensemble des $x$) pour lesquels $n_x\neq0$ soit localement fini;
ce sous-groupe contient le groupe libre $\bb{Z}^{(X^{(p)})}$ et lui est identique lorsque $X$ est noethérien.
On note $\mathfrak{Z}^{p+}(X)$ l'ensemble des éléments $\geq0$ de $\mathfrak{Z}^p(X)$.
Il est clair que le groupe ordonné $\mathfrak{Z}(X)$ contient la somme directe des sous-groupes ordonnés $\mathfrak{Z}^p(X)$, et est identique à cette somme directe lorsque $X$ est noethérien;
dans ce dernier cas, on considère $\mathfrak{Z}(X)$ comme \emph{gradué} par les $\mathfrak{Z}^p(X)$ pour $p\in\bb{N}$.
\end{env}

\begin{env}[21.6.3]
\label{IV.21.6.3}
Soient $Z=\sum_{x\in X}n_x\cdot\overline{\{x\}}$ un cycle sur $X$, $U$ un ouvert de $X$;
on appelle \emph{restriction} de $Z$ à $U$ et on note $Z|U$ le cycle $\sum_{x\in U}n_x\cdot(U\cap\overline{\{x\}})$ sur $U$;
on a $\Supp(Z|U)=\Supp(Z)\cap U$.
Il est clair que $Z\mapsto Z|U$ est un homomorphisme de groupes ordonnés de $\mathfrak{Z}(X)$ dans $\mathfrak{Z}(U)$ (resp.\ de $\mathfrak{Z}^p(X)$ dans $\mathfrak{Z}^p(U)$), et que si $V$ est un ouvert contenu dans $U$, on a $Z|V=(Z|U)|V$;
l'application $U\mapsto\mathfrak{Z}(U)$ (resp.\ $U\mapsto\mathfrak{Z}^p(U)$) est donc un \emph{préfaisceau} sur $X$ de groupes ordonnés commutatifs.
En fait ce préfaisceau est un \emph{faisceau}, somme directe des faisceaux $(i_x)_*(\bb{Z}_{\overline{\{x\}}})$, où $x$ parcourt $X$ (resp.\ $X^{(p)}$) et pour chaque $x\in X$, $i_x$ est l'injection canonique $\overline{\{x\}}\to X$ et $\bb{Z}_{\overline{\{x\}}}$ est le faisceau simple associé au faisceau constant $\bb{Z}$ sur l'espace $\overline{\{x\}}$:
cela résulte aussitôt de la description des sections d'une somme directe de faisceaux de groupes commutatifs (G, II, 2.7).
On note ce faisceau $\mathscr{Z}_X$ (resp.\ $\mathscr{Z}_X^p$).
On note $\mathscr{Z}_X^+$ (resp.\ $\mathscr{Z}_X^{p+}$) le sous-faisceau
\oldpage[IV-4]{272}
de monoïdes $U\mapsto\mathfrak{Z}^+(U)$ (resp.\ $U\mapsto\mathfrak{Z}^{p+}(U)$) de $\mathscr{Z}_X$ (resp.\ $\mathscr{Z}_X^p$).
Le faisceau $\mathscr{Z}_X$ est évidemment \emph{somme directe} des faisceaux $\mathscr{Z}_X^p$.

Notons enfin que les faisceaux $\mathscr{Z}_X^p$ (donc aussi $\mathscr{Z}_X$) sont \emph{flasques}.
Supposons en effet donnée une section $Z$ de $\mathscr{Z}_X^p$ au-dessus d'un ouvert $U$, de sorte que l'ensemble $S$ des $x\in U$ tels que $\operatorname{mult}_x(Z)\neq0$ est localement fini \emph{dans} $U$;
cet ensemble est aussi localement fini \emph{dans} $X$, car pour tout $z\in X$ et tout voisinage ouvert noethérien $V$ de $z$, $U\cap V$ est noethérien, donc ne contient qu'un nombre fini de points de $S$.
On définit donc une section $Z'$ de $\mathscr{Z}_X^p$ au-dessus de $X$ en posant $Z'=\sum_{x\in U}\operatorname{mult}_x(Z)\cdot\overline{\{x\}}$ et puisque l'adhérence de $x$ dans $U$ est $\overline{\{x\}}\cap U$, on a bien $Z'|U=Z$, d'où notre assertion.
\end{env}

\begin{env}[21.6.4]
\label{IV.21.6.4}
Nous nous proposons de définir un homomorphisme canonique
\[
\label{IV.21.6.4.1}
  c:\sh{D}iv_X\to\mathscr{Z}_X^1
  \tag{21.6.4.1}
\]
de faisceaux de groupes commutatifs.
Il suffit évidemment de définir un homomorphisme de $\sh{M}_X^*$ dans $\mathscr{Z}_X^1$, qui s'annule dans $\sh{O}_X^*$ \sref{IV.21.1.2};
or, par définition, $\sh{M}_X^*$ est le \emph{symétrisé} du faisceau de monoïdes $\sh{O}_X\cap\sh{M}_X^*$, et un homomorphisme $\sh{M}_X^*\to\mathscr{Z}_X^1$ est déterminé de façon unique par la donnée de sa restriction $\sh{S}(\sh{O}_X)\to\mathscr{Z}_X^1$ qui est un homomorphisme de faisceaux de \emph{monoïdes} soumis à la seule condition d'être nul dans $\sh{O}_X^*$;
il revient au même de dire que pour définir $c$, il suffit de définir un homomorphisme de \emph{faisceaux de monoïdes}
\[
\label{IV.21.6.4.2}
  c:\sh{D}iv_X^+\to\mathscr{Z}_X^{1+}.
  \tag{21.6.4.2}
\]
\end{env}

\begin{env}[21.6.5]
\label{IV.21.6.5}
Or, on a vu qu'à tout diviseur positif $D$ sur $X$ on fait correspondre le sous-préschéma fermé $Y(D)$ de $X$, défini par l'Idéal $\sh{J}_X(D)\subset\sh{O}_X$, et qui est régulièrement immergé et de codimension $1$.
En chacun des points maximaux $x$ de $Y(D)$, on a donc (\sref{IV.19.1.4} et \sref{IV.5.1.2}) $\codim(\overline{\{x\}},X)=\dim(\sh{O}_{X,x})=1$, autrement dit $x\in X^{(1)}$;
d'ailleurs l'ensemble de ces points maximaux est localement fini dans $X$ \sref{IV.3.1.6}, et $\sh{O}_{Y(D),x}$ est un anneau \emph{artinien}.
En tout point $x\in X^{(1)}$ qui n'est pas point maximal de $Y(D)$, on a nécessairement $x\notin Y(D)$, donc $\sh{O}_{Y(D),x}=0$.
Posons
\[
\label{IV.21.6.5.1}
  \operatorname{cyc}(D)=\sum_{x\in X^{(1)}}\operatorname{long}(\sh{O}_{Y(D),x})\cdot\overline{\{x\}}
  \tag{21.6.5.1}
\]
qui est donc un élément de $\mathfrak{Z}^1(X)$.
\end{env}

\begin{proposition}[21.6.6]
\label{IV.21.6.6}
L'application $D\mapsto\operatorname{cyc}(D)$ est un homomorphisme du monoïde $\operatorname{Div}^+(X)$ dans $\mathfrak{Z}^{1+}(X)$.
\end{proposition}

\begin{proof}
Tout revient à voir que pour deux diviseurs positifs $D$, $D'$, on a
\[
  \operatorname{cyc}(D+D')=\operatorname{cyc}(D)+\operatorname{cyc}(D').
\]
Or, on a \sref{IV.21.2.5} $\sh{J}_X(D+D')=\sh{J}_X(D)\sh{J}_X(D')$;
tout revient à voir que si $x\in X^{(1)}$, si l'on pose $A=\sh{O}_{X,x}$, et si $t$ et $t'$ sont deux éléments réguliers de $A$, on a $\operatorname{long}(A/tt'A)=\operatorname{long}(A/tA)+\operatorname{long}(A/t'A)$;
mais puisque $t$ est régulier, $tA/tt'A$ est isomorphe à $A/t'A$, d'où aussitôt la proposition.
\end{proof}

Il résulte aussitôt des définitions que pour tout ouvert $U\subset X$, on a
\[
  Y(D|U)=Y(D)\cap U,
\]
\oldpage[IV-4]{273}
donc $\operatorname{cyc}(D|U)=\operatorname{cyc}(D)|U$, et les homomorphismes $\operatorname{cyc}:\operatorname{Div}^+(U)\to\mathfrak{Z}^1(U)$ définissent donc un homomorphisme de faisceaux de monoïdes \sref{IV.21.6.4.2}, d'où l'homomorphisme cherché \sref{IV.21.6.4.1} de faisceaux de groupes.

On a $\Supp(\operatorname{cyc}(D))\subset\Supp(D)$ pour tout diviseur $D$ et
\begin{env}[21.6.6.2]
\label{IV.21.6.6.2}
\[
  \Supp(\operatorname{cyc}(D))=\Supp(D)\quad\text{lorsque}\quad D\geq0;
\]
\end{env}
on a déjà vu en effet la seconde relation \sref{IV.21.2.12};
lorsque $D$ est quelconque, la relation $D_x=0$ entraîne l'existence d'un voisinage ouvert $U$ de $x$ tel que $D|U=0$, d'où $\operatorname{cyc}(D)|U=0$, ce qui prouve notre assertion.

\begin{env}[21.6.7]
\label{IV.21.6.7}
L'homomorphisme \sref{IV.21.6.4.1} donne, par passage aux groupes de sections au-dessus de $X$, un \emph{homomorphisme croissant de groupes ordonnés} $\operatorname{Div}(X)\to\mathfrak{Z}^1(X)$, que nous noterons encore $D\mapsto\operatorname{cyc}(D)$;
le nombre $\operatorname{mult}_x(\operatorname{cyc}(D))$ pour $x\in X^{(1)}$ se note encore $\operatorname{mult}_x(D)$ ou $\operatorname{mult}_{\overline{\{x\}}}(D)$ et s'appelle \emph{multiplicité du diviseur $D$ au point $x$}, ou \emph{multiplicité du cycle premier $\overline{\{x\}}$ dans $D$};
c'est un entier positif ou négatif, égal comme on l'a vu à $\operatorname{long}(\sh{O}_{Y(D),x})$ lorsque $D$ est un diviseur \emph{positif};
on a donc par définition
\[
\label{IV.21.6.7.1}
  \operatorname{cyc}(D)=\sum_{x\in X^{(1)}}\operatorname{mult}_x(D)\cdot\overline{\{x\}},
  \tag{21.6.7.1}
\]
et en vertu du fait que $D\mapsto\operatorname{cyc}(D)$ est un homomorphisme, on a
\[
\label{IV.21.6.7.2}
  \operatorname{mult}_x(-D)=-\operatorname{mult}_x(D),\qquad
  \operatorname{mult}_x(D+D')=\operatorname{mult}_x(D)+\operatorname{mult}_x(D')
  \tag{21.6.7.2}
\]
pour deux diviseurs quelconques $D$, $D'$.

Pour toute fonction méromorphe régulière $f$ sur $X$, on pose en particulier, pour $x\in X^{(1)}$
\[
\label{IV.21.6.7.3}
  o_x(f)=\operatorname{mult}_x(\operatorname{div}(f))
  \tag{21.6.7.3}
\]
et on dit que $o_x(f)$ (entier positif ou négatif) est \emph{l'ordre de $f$ au point} $x\in X^{(1)}$.
Si $f_x\in\sh{O}_{X,x}$ (élément régulier de cet anneau local par hypothèse), on a donc
\[
\label{IV.21.6.7.4}
  o_x(f)=\operatorname{long}(\sh{O}_{X,x}/f_x\sh{O}_{X,x});
  \tag{21.6.7.4}
\]
pour deux fonctions méromorphes régulières $f$, $f'$ sur $X$, on a
\[
\label{IV.21.6.7.5}
  o_x(ff')=o_x(f)+o_x(f'),\qquad o_x(f^{-1})=-o_x(f)
  \tag{21.6.7.5}
\]
pour tout $x\in X^{(1)}$.
Les cycles $1$-codimensionnels
\[
  Z^+(f)=\sum_{x\in X^{(1)}}(o_x(f))^+\cdot\overline{\{x\}},\qquad
  Z^-(f)=\sum_{x\in X^{(1)}}(o_x(f))^-\cdot\overline{\{x\}}
\]
sont appelés respectivement le \emph{cycle des zéros} et le \emph{cycle des pôles} (ou \emph{cycle polaire}) de $f$, et l'on a évidemment $\operatorname{cyc}(\operatorname{div}(f))=Z^+(f)-Z^-(f)$.
Les cycles $1$-codimensionnels de la forme $\operatorname{cyc}(\operatorname{div}(f))$ sont appelés \emph{principaux} (ou \emph{linéairement équivalents à zéro}) et forment un sous-groupe $\mathfrak{Z}^1_\mathrm{princ}(X)$ de $\mathfrak{Z}^1(X)$.
Les sections au-dessus de $X$ de l'image $c(\sh{D}iv_X)\subset\mathscr{Z}_X^1$ sont
\oldpage[IV-4]{274}
appelées \emph{cycles localement principaux};
un tel cycle $Z$ est donc caractérisé par le fait que, pour tout $y\in X$, il y a un voisinage ouvert $U$ de $y$ dans $X$ et une fonction méromorphe régulière $f$ sur $U$ telle que $Z|U=\sum_{x\in U\cap X^{(1)}}o_x(f)(\overline{\{x\}}\cap U)$.
Si $Z=\sum_{x\in X^{(1)}}n_x\cdot\overline{\{x\}}$, il revient au même de dire que, pour tout $y\in X$, si l'on pose $T_y=\Spec(\sh{O}_{X,y})$ et $Z_y=\sum_{x\in X^{(1)}\cap T_y}n_x(\overline{\{x\}}\cap T_y)$, $Z_y$ est \emph{principal}, autrement dit il existe une fonction méromorphe régulière $g$ sur $T_y$ telle que $Z_y=\operatorname{cyc}(\operatorname{div}(g))$.
Cela résulte aussitôt en effet de la définition précédente et de \sref{IV.20.3.7} qui établit une correspondance biunivoque entre les fonctions méromorphes régulières sur $T_y$ et les germes de fonctions méromorphes régulières sur les voisinages ouverts de $y$ dans $X$ lorsque $X$ est localement noethérien.
On exprime encore le fait que $Z_y$ est principal en disant que $Z$ est \emph{principal au point $y$}.
L'ensemble des points où $Z$ est principal est évidemment \emph{ouvert}, en vertu de ce qui précède.

On déduit de l'homomorphisme canonique $\operatorname{cyc}:\operatorname{Div}(X)\to\mathfrak{Z}^1(X)$ un homomorphisme $\operatorname{Div}(X)/\operatorname{Div.princ}(X)\to\mathfrak{Z}^1(X)/\mathfrak{Z}^1_\mathrm{princ}(X)$, en vertu de la définition de $\mathfrak{Z}^1_\mathrm{princ}(X)$.
On dit que le groupe $\mathfrak{Z}^1(X)/\mathfrak{Z}^1_\mathrm{princ}(X)$ est le \emph{groupe des classes de cycles $1$-codimensionnels} de $X$ et on le note $\mathfrak{Cl}(X)$.
Deux éléments de $\mathfrak{Z}^1(X)$ ayant même image dans $\mathfrak{Cl}(X)$ sont appelés \emph{linéairement équivalents}.
\end{env}

\begin{env}[21.6.8]
\label{IV.21.6.8}
Considérons en particulier le cas où $X=\Spec(A)$, où $A$ est un anneau \emph{noethérien intègre et intégralement clos}.
Alors $X^{(1)}$ est l'ensemble des \emph{idéaux premiers de hauteur $1$} de $A$, et $\mathfrak{Z}^1(X)$ s'identifie donc au \emph{groupe des diviseurs} (au sens de N.~Bourbaki) de l'anneau de Krull $A$ (Bourbaki, \emph{Alg.\ comm.}, chap.~VII, \textsection~1, n\textsuperscript{o}~3, cor.\ du th.~2 et n\textsuperscript{o}~6, th.~3).

Comme d'autre part, les fonctions méromorphes régulières sur $X$ s'identifient alors aux éléments non nuls du corps des fractions $K$ de $A$, l'application $f\mapsto\operatorname{cyc}(\operatorname{div}(f))$ de $M(X)$ dans $\mathfrak{Z}^1(X)$ s'identifie à l'application notée $f\mapsto\operatorname{div}(f)$ dans Bourbaki (\emph{loc.\ cit.}, \textsection~1, n\textsuperscript{o}~1);
$\mathfrak{Z}^1_\mathrm{princ}(X)$ s'identifie donc au \emph{groupe des diviseurs principaux} de $A$ au sens de Bourbaki, et $\mathfrak{Cl}(X)$ au \emph{groupe des classes de diviseurs} de $A$ au sens de Bourbaki (\emph{loc.\ cit.}, \textsection~1, n\textsuperscript{o}~2 et n\textsuperscript{o}~10).
\end{env}

\begin{theorem}[21.6.9]
\label{IV.21.6.9}
Soit $X$ un préschéma localement noethérien normal.
\begin{enumerate}[label=(\roman*)]
  \item L'homomorphisme canonique $\operatorname{cyc}:\operatorname{Div}(X)\to\mathfrak{Z}^1(X)$ est injectif et son image est formée des cycles localement principaux.
  \item Les conditions suivantes sont équivalentes:
  \begin{enumerate}[label=\emph{\alph*)}]
    \item L'homomorphisme $\operatorname{cyc}:\operatorname{Div}(X)\to\mathfrak{Z}^1(X)$ est bijectif.
    \item Tout cycle $1$-codimensionnel est localement principal.
    \item Pour tout $x\in X$, l'anneau local $\sh{O}_{X,x}$ est factoriel.
  \end{enumerate}
  (On dit alors que $X$ est un préschéma \emph{localement factoriel}.)
\end{enumerate}
\end{theorem}

\begin{proof}
\begin{enumerate}[label=(\roman*)]
  \item Il suffit de démontrer que
\[
\label{IV.21.6.9.1}
  \operatorname{cyc}^{-1}(\mathfrak{Z}^{1+}(X))=\operatorname{Div}^+(X),
  \tag{21.6.9.1}
\]
puisque l'on a $\operatorname{Div}^+(X)\cap(-\operatorname{Div}^+(X))=0$ et $\mathfrak{Z}^{1+}(X)\cap(-\mathfrak{Z}^{1+}(X))=0$.
On est donc ramené à prouver que si $D$ est un diviseur sur $X$ tel que $\operatorname{mult}_x(D)\geq0$ pour tout $x\in X^{(1)}$, on a $D\geq0$.
Or, pour tout $x\notin X^{(1)}$, l'anneau local $\sh{O}_{X,x}$ est intègre et intégralement clos et
\oldpage[IV-4]{275}
de dimension $0$ ou $\geq2$, donc on a $\operatorname{prof}(\sh{O}_{X,x})=0$ ou $\operatorname{prof}(\sh{O}_{X,x})\geq2$ \sref[0]{0.16.5.1}.
D'autre part, aux points $x\in X^{(1)}$, l'anneau $\sh{O}_{X,x}$ est un anneau de valuation discrète (\sref[II]{II.7.1.6}), donc $\operatorname{prof}(\sh{O}_{X,x})=1$ \sref[0]{0.16.5.1};
les seuls points de $X$ tels que $\operatorname{prof}(\sh{O}_{X,x})=1$ sont donc les points de $X^{(1)}$, et la conclusion résulte par suite de \sref{IV.21.1.8}.

  \item L'équivalence de \emph{a)} et \emph{b)} est claire en vertu de (i).
D'après la caractérisation des cycles localement principaux donnée dans \sref{IV.21.6.7}, et la relation entre cycles $1$-codimensionnels sur $\Spec(A)$ et diviseurs (au sens de Bourbaki) de l'anneau $A$ lorsque $A$ est un anneau noethérien intégralement clos \sref{IV.21.6.8}, la condition \emph{b)} équivaut à dire que pour tout $x\in X$, tout diviseur de l'anneau $\sh{O}_{X,x}$ est principal, autrement dit que l'anneau $\sh{O}_{X,x}$ est factoriel (Bourbaki, \emph{Alg.\ comm.}, chap.~VII, \textsection~3, n\textsuperscript{o}~1), d'où l'équivalence de \emph{b)} et \emph{c)}.
\end{enumerate}
\end{proof}

\begin{env}[21.6.9.2]
\label{IV.21.6.9.2}
Lorsque $A$ est un anneau noethérien \emph{factoriel}, il est clair que $X=\Spec(A)$ est localement factoriel (Bourbaki, \emph{Alg.\ comm.}, chap.~VII, \textsection~3, n\textsuperscript{o}~4, prop.~3).
Si, dans ce cas, on écrit un élément $f\neq0$ du corps des fractions $K$ de $A$ sous la forme $r/s$, où $r$ et $s$ sont deux éléments \emph{étrangers} de $A$, dont les diviseurs sont bien déterminés (\emph{loc.\ cit.}, \textsection~3, n\textsuperscript{o}~3), ces diviseurs s'identifient respectivement au \emph{diviseur des zéros} et au \emph{diviseur des pôles} de $f$ \sref{IV.21.6.7}.
\end{env}

\begin{corollary}[21.6.10]
\label{IV.21.6.10}
Soit $X$ un préschéma localement noethérien normal.
\begin{enumerate}[label=(\roman*)]
  \item Il existe un homomorphisme canonique injectif
\[
\label{IV.21.6.10.1}
  \operatorname{Pic}(X)\to\mathfrak{Cl}(X).
  \tag{21.6.10.1}
\]

  \item Si $X$ est localement factoriel, l'homomorphisme \sref{IV.21.6.10.1} est bijectif, et réciproquement.
\end{enumerate}
\end{corollary}

\begin{proof}
On a vu \sref{IV.21.6.7} que l'image de $\operatorname{Div.princ}(X)$ par l'homomorphisme $\operatorname{cyc}:\operatorname{Div}(X)\to\mathfrak{Z}^1(X)$ est $\mathfrak{Z}^1_\mathrm{princ}(X)$;
on déduit donc de \sref{IV.21.6.9} que l'homomorphisme $\operatorname{Div}(X)/\operatorname{Div.princ}(X)\to\mathfrak{Z}^1(X)/\mathfrak{Z}^1_\mathrm{princ}(X)=\mathfrak{Cl}(X)$ déduit de $\operatorname{cyc}$ est injectif, et qu'il est bijectif si et seulement si $X$ est localement factoriel.
La conclusion résulte donc de ce que, lorsque $X$ est localement noethérien et réduit, l'homomorphisme canonique $\operatorname{Div}(X)/\operatorname{Div.princ}(X)\to\operatorname{Pic}(X)$ \sref{IV.21.3.3.1} est bijectif \sref{IV.21.3.4}, (ii).
\end{proof}

\begin{corollary}[21.6.11]
\label{IV.21.6.11}
Soit $X$ un préschéma localement noethérien et localement factoriel.
Alors le faisceau $\sh{D}iv_X$ est flasque, et pour tout ouvert $U$ de $X$, l'homomorphisme canonique $\operatorname{Pic}(X)\to\operatorname{Pic}(U)$ est surjectif.
\end{corollary}

\begin{proof}
Tenant compte de \sref{IV.21.6.9}, (ii), la première assertion équivaut à dire que le faisceau $\mathscr{Z}_X^1$ est flasque, ce qui a été vu plus haut \sref{IV.21.6.3}.
Pour tout ouvert $U$ de $X$, l'homomorphisme canonique $\mathfrak{Z}^1(X)\to\mathfrak{Z}^1(U)$ est donc surjectif;
comme en vertu de \sref{IV.21.6.10}, l'homomorphisme $\operatorname{Pic}(X)\to\operatorname{Pic}(U)$ s'identifie canoniquement à $\mathfrak{Z}^1(X)/\mathfrak{Z}^1_\mathrm{princ}(X)\to\mathfrak{Z}^1(U)/\mathfrak{Z}^1_\mathrm{princ}(U)$, il est aussi surjectif.
\end{proof}

\begin{proposition}[21.6.12]
\label{IV.21.6.12}
Soit $X$ un préschéma noethérien réduit.
Soit $(U_\lambda)_{\lambda\in L}$ une famille filtrante décroissante d'ouverts de $X$ vérifiant les conditions suivantes:
\begin{enumerate}
  \item[\rm 1\textsuperscript{o}] Si $Y_\lambda=X-U_\lambda$, on a $\codim(Y_\lambda,X)\geq2$ pour tout $\lambda\in L$.
  \item[\rm 2\textsuperscript{o}] Pour tout $x\in\bigcap_{\lambda\in L}U_\lambda$, l'anneau $\sh{O}_{X,x}$ est factoriel.
\end{enumerate}
Alors il existe des isomorphismes canoniques
\[
\label{IV.21.6.12.1}
  \varinjlim\operatorname{Div}(U_\lambda)\simto\mathfrak{Z}^1(X),\qquad
  \varinjlim\operatorname{Pic}(U_\lambda)\simto\mathfrak{Cl}(X)
  \tag{21.6.12.1}
\]
tels que, pour tout ouvert $V$ de $X$, les diagrammes
\[
\label{IV.21.6.12.2}
  \xymatrix{
    \varinjlim\operatorname{Div}(U_\lambda) \ar[r]^-\sim \ar[d] & \mathfrak{Z}^1(X) \ar[d] & & \varinjlim\operatorname{Pic}(U_\lambda) \ar[r]^-\sim \ar[d] & \mathfrak{Cl}(X) \ar[d]\\
    \varinjlim\operatorname{Div}(U_\lambda\cap V) \ar[r]^-\sim & \mathfrak{Z}^1(V) & & \varinjlim\operatorname{Pic}(U_\lambda\cap V) \ar[r]^-\sim & \mathfrak{Cl}(V)
  }
  \tag{21.6.12.2}
\]
soient commutatifs.
\end{proposition}

\begin{proof}
L'hypothèse 1\textsuperscript{o} sur $U_\lambda$ implique que pour tout $\lambda\in L$, on a $X^{(1)}\subset U_\lambda$ \sref{IV.5.1.3.1} donc l'homomorphisme de restriction $\mathfrak{Z}^1(X)\to\mathfrak{Z}^1(U_\lambda)$ est bijectif et par suite on a un isomorphisme canonique $\varinjlim\mathfrak{Z}^1(U_\lambda)\simto\mathfrak{Z}^1(X)$.
Les homomorphismes canoniques $\operatorname{cyc}:\operatorname{Div}(U_\lambda)\to\mathfrak{Z}^1(U_\lambda)$ définissent donc, par passage à la limite inductive, le premier des homomorphismes canoniques \sref{IV.21.6.12.1}, et le second s'en déduit par passage aux quotients.
En outre, il résulte de la condition 1\textsuperscript{o} que les $U_\lambda$ sont denses dans $X$, donc schématiquement denses puisque $X$ est réduit \sref{IV.11.10.4}, et par suite toute fonction méromorphe sur $X$ est entièrement déterminée par sa restriction à chaque $U_\lambda$;
on déduit aussitôt de là que dans l'isomorphisme $\mathfrak{Z}^1(X)\simto\mathfrak{Z}^1(U_\lambda)$, l'image de $\mathfrak{Z}^1_\mathrm{princ}(X)$ est $\mathfrak{Z}^1_\mathrm{princ}(U_\lambda)$, donc on a aussi un isomorphisme canonique $\varinjlim\mathfrak{Cl}(U_\lambda)\simto\mathfrak{Cl}(X)$.
Le second des homomorphismes canoniques \sref{IV.21.6.12.1} sera donc un isomorphisme si le premier l'est, et il reste donc à prouver ce dernier point, la commutativité des diagrammes \sref{IV.21.6.12.2} étant triviale.

\oldpage[IV-4]{276}
Montrons d'abord que l'homomorphisme $\varinjlim\operatorname{Div}(U_\lambda)\to\mathfrak{Z}^1(X)$ est \emph{injectif}.
Posons $T=\bigcap_\lambda U_\lambda$, et notons que les $U_\lambda$ forment un \emph{système fondamental de voisinages} de $T$.
En effet, pour tout ouvert $V\supset T$, $X-V$ est un sous-espace fermé de $X$, donc un espace noethérien dont toute partie fermée irréductible admet un point générique;
comme les ensembles $(X-V)\cap(X-U_\lambda)$ sont fermés dans $X-V$, forment une famille filtrante croissante et ont pour réunion $X-V$, notre assertion résulte de \sref[0]{0.9.2.4}.
Cela étant, il s'agit de voir que si $D\in\operatorname{Div}(U_\lambda)$ est tel que $\operatorname{cyc}(D)=0$ dans $\mathfrak{Z}^1(U_\lambda)$, alors il existe $\mu\geq\lambda$ tel que $D|U_\mu=0$.
En vertu de ce qui précède, il suffira de voir que pour tout $x\in T$, on a $D_x=0$;
en effet, il y aura alors un voisinage $W(x)$ de $x$ dans $X$ tel que $D|W(x)=0$, et la réunion des $W(x)$ pour $x\in T$ contient un $U_\mu$.
Compte tenu de \sref{IV.21.4.6}, on est donc ramené au cas où $X=\Spec(\sh{O}_{X,x})$ avec $x\in T$;
mais par hypothèse $\sh{O}_{X,x}$ est factoriel, donc intègre et intégralement clos, et $\Spec(\sh{O}_{X,x})$ est alors normal;
donc la conclusion résulte de \sref{IV.21.6.9}, (i).

\oldpage[IV-4]{277}
Pour prouver que l'homomorphisme $\varinjlim\operatorname{Div}(U_\lambda)\to\mathfrak{Z}^1(X)$ est bijectif, il suffira de prouver de même que pour tout $Z\in\mathfrak{Z}^1(X)$ et tout $x\in T$, il y a un voisinage $W(x)$ de $x$ dans $X$ et un diviseur $D^{(x)}$ sur $W(x)$ tel que $\operatorname{cyc}(D^{(x)})=Z|W(x)$.
En effet, on pourra alors recouvrir l'ensemble quasi-compact $T$ par un nombre fini de $W(x_i)$, et en vertu de la première partie de la démonstration, puisque les couples $(i,j)$ sont en nombre fini, il existera un indice $\lambda$ tel que dans chacun des $W(x_i)\cap W(x_j)\cap U_\lambda$, les restrictions de $D^{(x_i)}$ et $D^{(x_j)}$ coïncident;
comme on peut supposer en outre que $U_\lambda$ est contenu dans la réunion des $W(x_i)$, on voit que les restrictions $D^{(x_i)}|(W(x_i)\cap U_\lambda)$ sont les restrictions d'un même diviseur $D\in\operatorname{Div}(U_\lambda)$ qui sera tel que $\operatorname{cyc}(D)=Z|U_\lambda$.
On est donc ramené encore au cas où $X=\Spec(\sh{O}_{X,x})$ avec $x\in T$, mais comme $\sh{O}_{X,x}$ est factoriel, il en est de même de ses localisés (Bourbaki, \emph{Alg.\ comm.}, chap.~VII, \textsection~3, n\textsuperscript{o}~4, prop.~3), et il suffit d'appliquer \sref{IV.21.6.9}, (ii).
\end{proof}

\begin{corollary}[21.6.13]
\label{IV.21.6.13}
Soit $A$ un anneau local noethérien intègre et intégralement clos et de dimension $\geq2$.
Posons $X=\Spec(A)$, et soient $a$ le point fermé de $X$, $U=X-\{a\}$.
Pour que $A$ soit factoriel, il faut et il suffit que $U$ soit localement factoriel et que $\operatorname{Pic}(U)=0$.
\end{corollary}

\begin{proof}
En effet, dire que $A$ est factoriel équivaut à dire que $\mathfrak{Cl}(X)=0$ (Bourbaki, \emph{Alg.\ comm.}, chap.~VII, \textsection~1, n\textsuperscript{o}~4, cor.\ du th.~2 et \textsection~3, n\textsuperscript{o}~2, th.~1);
il suffit donc d'utiliser l'existence du second isomorphisme \sref{IV.21.6.12.1}, en prenant la famille $(U_\lambda)$ restreinte au seul ouvert $U$.
\end{proof}

\begin{corollary}[21.6.14]
\label{IV.21.6.14}
Soient $A$ un anneau local noethérien de dimension $\geq2$, $X=\Spec(A)$, $a$ le point fermé de $X$, $U=X-\{a\}$.
Les conditions suivantes sont équivalentes:
\begin{enumerate}[label=\emph{\alph*)}]
  \item $A$ est factoriel.
  \item $\operatorname{Pic}(U)=0$ et $\operatorname{prof}(A)\geq2$ (conditions que nous exprimerons plus loin \sref{IV.21.13} en disant que l'anneau $A$ est \emph{parafactoriel}), et $U$ est localement factoriel.
\end{enumerate}
\end{corollary}

\begin{proof}
Il est clair que si $A$ est factoriel, il est normal, donc $\operatorname{prof}(A)\geq2$ puisque $\dim(A)\geq2$ \sref[0]{0.16.5.1}, et \sref{IV.21.6.13} montre que \emph{a)} implique \emph{b)}.
Inversement, si \emph{b)} est vérifiée, il suffit de voir que $A$ est intégralement clos pour en conclure par \sref{IV.21.6.13} que \emph{b)} entraîne \emph{a)}.
En vertu du critère de Serre \sref{IV.5.8.6}, il suffit de vérifier pour $X$ les conditions $(\mathrm{R}_1)$ et $(\mathrm{S}_2)$.
Or, $U$ étant localement factoriel vérifie ces conditions, et l'hypothèse $\operatorname{prof}(A)\geq2$ entraîne que $X$ les vérifie également.
\end{proof}
\subsection{Interprétation des cycles positifs 1-codimensionnels en termes de sous-préschémas}
\label{subsection:IV.21.7}
\label{IV.21.7}

\begin{env}[21.7.1]
\label{IV.21.7.1}
Soient $X$ un préschéma localement noethérien, $C=\sum_{x\in X^{(1)}}n_x\cdot\overline{\{x\}}$ un cycle $1$-codimensionnel positif (de sorte que l'on a $n_x\geq0$ pour tout $x\in X^{(1)}$, et $n_x=0$ sauf en un ensemble localement fini de points).
Désignons par $Y(C)$ le sous-préschéma fermé de $X$, \emph{image fermée} (\sref[I]{I.9.5.3} et \sref{IV.9.5.1}) du préschéma \emph{somme} $Y'(C)=\coprod_{x\in X^{(1)}}\Spec(\sh{O}_{X,x}/\mathfrak{m}_x^{n_x})$ par le morphisme canonique, et par $\sh{J}_X(C)$ (ou $\sh{J}(C)$) l'Idéal de $\sh{O}_X$ définissant $Y(C)$.
\end{env}

\begin{proposition}[21.7.2]
\label{IV.21.7.2}
Soit $X$ un préschéma localement noethérien vérifiant la condition $(\mathrm{R}_1)$ \sref{IV.5.8.2}.
Pour qu'un sous-préschéma fermé $Y$ de $X$ soit de la forme $Y(C)$, où $C$ est un cycle $1$-codimensionnel positif, il faut et il suffit qu'il vérifie les deux conditions suivantes:
\begin{enumerate}[label=(\roman*)]
  \item $Y$ est purement de codimension $1$.
\oldpage[IV-4]{278}
  \item $Y$ vérifie la condition $(\mathrm{S}_1)$, autrement dit \sref{IV.5.7.5} n'a pas de cycle premier associé immergé.
\end{enumerate}
On a alors
\[
\label{IV.21.7.2.1}
  \operatorname{mult}_x(C)=\operatorname{long}\sh{O}_{Y,x}.
  \tag{21.7.2.1}
\]
L'application $C\mapsto Y(C)$ est une bijection de $\mathfrak{Z}^{1+}(X)$ sur l'ensemble des sous-préschémas fermés de $X$ vérifiant les conditions (i) et (ii).
\end{proposition}

\begin{proof}
Les conditions sont nécessaires (sans supposer que $X$ vérifie $(\mathrm{R}_1)$).
En effet, la question étant locale sur $X$, on peut supposer que $X=\Spec(A)$, où $A$ est un anneau noethérien;
on a alors $Y(C)=\Spec(A/\mathfrak{J})$, où $\mathfrak{J}$ est, par définition \sref[I]{I.9.5.1}, le noyau de l'homomorphisme canonique $A\to\bigoplus_i A_{\mathfrak{p}_i}/\mathfrak{p}_i^{n_i}A_{\mathfrak{p}_i}$, où les $\mathfrak{p}_i$ sont les idéaux premiers correspondant aux points $x\in X^{(1)}$ pour lesquels $n_x>0$, et où l'on a posé $n_i=n_x$.
L'image réciproque $\mathfrak{q}_i$ de $\mathfrak{p}_i^{n_i}A_{\mathfrak{p}_i}$ dans $A$ est un idéal $\mathfrak{p}_i$-primaire et on a $\mathfrak{J}=\bigcap_i\mathfrak{q}_i$.
D'ailleurs, comme les $x\in X^{(1)}$ sont tels que $\overline{\{x\}}$ est de codimension $1$, aucun point de $X^{(1)}$ ne peut être adhérent à un autre point de $X^{(1)}$.
Donc les $\mathfrak{p}_i$ sont les idéaux premiers minimaux de $A/\mathfrak{J}$ et l'ensemble des $\mathfrak{p}_i$ est égal à $\operatorname{Ass}(A/\mathfrak{J})$, ce qui prouve les conditions (i) et (ii).

Les conditions sont suffisantes.
Désignant par $\sh{J}$ l'Idéal de $\sh{O}_X$ définissant $Y$, l'hypothèse (ii) implique que $\operatorname{Ass}(\sh{O}_X/\sh{J})$ est identique à l'ensemble $F$ des points maximaux de $Y$, et l'hypothèse (i) implique que $F$ est contenu dans $X^{(1)}$;
donc \sref{IV.3.2.6} $\sh{O}_X/\sh{J}$ s'identifie à un sous-$\sh{O}_X$-Module de $\bigoplus_{x\in F}\mathscr{G}(x)$, où pour chaque $x\in F$, $\mathscr{G}(x)$ est un $\sh{O}_X$-Module irredondant tel que $\operatorname{Ass}(\mathscr{G}(x))=\{x\}$.
Or, puisque $X$ vérifie $(\mathrm{R}_1)$, pour chaque $x\in F$, $\sh{O}_{X,x}$ est un anneau de valuation discrète et par suite les idéaux primaires non nuls de cet anneau sont les puissances $\mathfrak{m}_x^{n_x}$ de l'idéal maximal;
supposant encore que $X=\Spec(A)$, on en conclut que $\sh{J}=\widetilde{\mathfrak{J}}$, où $\widetilde{\mathfrak{J}}=\bigcap_i\mathfrak{q}_i$, les $\mathfrak{q}_i$ étant les images réciproques dans $A$ d'idéaux $\mathfrak{p}_i^{n_i}A_{\mathfrak{p}_i}$, où les $\mathfrak{p}_i$ correspondent aux points de $F$.
On voit donc bien que $Y$ est de la forme $Y(C)$.
\end{proof}

\begin{corollary}[21.7.3]
\label{IV.21.7.3}
Soit $X$ un préschéma localement noethérien.
Pour tout diviseur positif $D$ sur $X$, tel que $X$ vérifie $(\mathrm{R}_1)$ aux points maximaux du sous-préschéma fermé $Y(D)$ \sref{IV.21.2.12}, $Y(D)$ majore le sous-préschéma fermé $Y(\operatorname{cyc}(D))$ \sref{IV.21.7.1}, et a même espace sous-jacent;
pour que ces deux sous-préschémas soient égaux, il faut et il suffit que $Y(D)$ vérifie la condition $(\mathrm{S}_1)$, ou encore que, pour tout $z\in Y(D)$ distinct d'un point maximal de $Y(D)$, on ait $\operatorname{prof}(\sh{O}_{X,z})\geq2$ (condition toujours vérifiée lorsque $X$ est \emph{normal}).
\end{corollary}

\begin{proof}
En effet, la question étant locale, on peut toujours supposer que $\sh{J}_X(D)=t\sh{O}_X$, $t$ étant une section régulière de $\sh{O}_X$ au-dessus de $X$;
en tout point maximal $x$ de $Y(D)$, appartenant nécessairement à $X^{(1)}$, $\sh{O}_{X,x}$ est par hypothèse un anneau de valuation discrète, donc $t_x\sh{O}_{X,x}=\mathfrak{m}_x^{n_x}$ où $n_x=\operatorname{mult}_x(D)$.
On peut supposer $X=\Spec(A)$, et alors, si $\sh{J}_X(\operatorname{cyc}(D))=\widetilde{\mathfrak{J}}$, il résulte de ce qui précède et de \sref{IV.21.7.1} que $\mathfrak{J}=tA$ et $\widetilde{\mathfrak{J}}$ ont mêmes \emph{idéaux primaires non immergés}, donc $\widetilde{\mathfrak{J}}\subset\mathfrak{J}$, puisque $\widetilde{\mathfrak{J}}$ est l'intersection
\oldpage[IV-4]{279}
de ces idéaux primaires \sref{IV.21.7.2}.
Cela prouve que $Y(D)$ majore $Y(\operatorname{cyc}(D))$ et que ces deux sous-préschémas sont égaux si et seulement si $Y(D)$ n'a pas de cycle premier associé immergé (autrement dit vérifie $(\mathrm{S}_1)$).
Comme pour tout $x\in Y(D)$, il y a par hypothèse un voisinage ouvert $U$ de $x$ dans $X$ et un élément régulier $t\in\Gamma(U,\sh{O}_X)$ tel que $\sh{O}_{Y(D)}|(U\cap Y(D))$ soit restriction à $Y(D)\cap U$ de $\sh{O}_U/t\sh{O}_U$, dire que $Y(D)$ vérifie $(\mathrm{S}_1)$ signifie aussi qu'en tout point $z\in Y(D)$ non maximal, on a $\operatorname{prof}(\sh{O}_{X,z}/t_z\sh{O}_{X,z})\geq1$, c'est-à-dire \sref[0]{0.16.4.6} $\operatorname{prof}(\sh{O}_{X,z})\geq2$.
L'assertion concernant le cas où $X$ est normal est alors immédiate puisque $X$ vérifie $(\mathrm{S}_2)$ et qu'aux points non maximaux $z$ de $Y(D)$ on a $\dim(\sh{O}_{X,z})\geq2$ \sref[0]{0.16.3.4}.
\end{proof}

\begin{env}[21.7.3.1]
\label{IV.21.7.3.1}
On voit donc que lorsque $X$ est \emph{normal}, $\operatorname{Div}^+(X)$ s'identifie canoniquement à l'ensemble des sous-préschémas fermés $Y$ de $X$ vérifiant les conditions (i) et (ii) de \sref{IV.21.7.2} et \emph{régulièrement immergés} dans $X$.
\end{env}

\begin{proposition}[21.7.4]
\label{IV.21.7.4}
Soit $X$ un préschéma localement noethérien, réduit en chacun de ses points isolés.
Les conditions suivantes sont équivalentes:
\begin{enumerate}
  \item[\emph{a)}] L'homomorphisme canonique $c:\sh{D}iv_X\to\mathscr{Z}_X^1$ \sref{IV.21.6.4} est un isomorphisme de faisceaux de groupes ordonnés.
  \item[\emph{a$'$)}] Tout cycle premier $1$-codimensionnel de $X$ est l'image par $\operatorname{cyc}$ d'un diviseur positif, et l'homomorphisme canonique $c:\sh{D}iv_X\to\mathscr{Z}_X^1$ est injectif.
  \item[\emph{a$''$)}] Pour tout sous-préschéma fermé intègre $Y$ de $X$, de codimension $1$, l'immersion canonique $Y\to X$ est régulière.
  \item[\emph{b)}] $X$ est normal et l'homomorphisme $\operatorname{cyc}:\operatorname{Div}(X)\to\mathfrak{Z}^1(X)$ est bijectif.
  \item[\emph{c)}] $X$ est localement factoriel.
\end{enumerate}
\end{proposition}

\begin{proof}
L'équivalence de \emph{b)} et \emph{c)} a été démontrée dans \sref{IV.21.6.9}, ainsi que le fait que \emph{c)} entraîne \emph{a)}.
De plus \emph{b)} entraîne \emph{a$''$)} en vertu de \sref{IV.21.7.3.1}, et \emph{a)} implique trivialement \emph{a$'$)}.
Il reste à voir que \emph{a$'$)} ou \emph{a$''$)} entraîne \emph{c)}.

Supposons donc vérifiée la condition \emph{a$'$)}, et montrons d'abord que $X$ est normal.
Notons en premier lieu que si $X$ vérifie \emph{a$'$)}, il en est de même de tout schéma local $\Spec(\sh{O}_{X,x})$.
Considérons alors $x\in X^{(1)}$, de sorte que $A=\sh{O}_{X,x}$ est un anneau local noethérien de dimension $1$.
Appliquant la condition \emph{a$'$)} à $\Spec(A)$ et au cycle premier $1$-codimensionnel formé du point fermé de $\Spec(A)$, on voit que dans $A$ l'idéal maximal est engendré par un seul élément régulier dans $A$;
donc \sref[0]{0.17.1.1} $A$ est un anneau régulier.
Comme les anneaux localisés $A_{\mathfrak p}$ sont aussi réguliers \sref[0]{0.17.3.2}, on voit que $X$ est régulier en tous ses points maximaux non isolés;
comme il est aussi réduit (donc régulier) en ses points isolés par hypothèse, on en conclut que $X$ vérifie la condition $(\mathrm{R}_1)$.
Montrons en outre que $X$ vérifie aussi $(\mathrm{S}_1)$, autrement dit que pour tout $x\in X$ tel que $\dim(\sh{O}_{X,x})\geq1$, on a $\mathfrak{m}_x\notin\operatorname{Ass}(\sh{O}_{X,x})$ \sref[0]{0.16.4.6}.
En effet, l'hypothèse \emph{a$'$)} appliquée à $X_1=\Spec(\sh{O}_{X,x})$ montre qu'il existe sur ce préschéma au moins un diviseur non nul, autrement dit que l'on a $\sh{M}_{X_1}^*\neq\sh{O}_{X_1}^*$, et il suffit d'appliquer \sref{IV.21.1.10}.
On déduit donc déjà de ces résultats que $X$ est réduit \sref{IV.5.8.5}.
Prouvons ensuite qu'il vérifie la condition $(\mathrm{S}_2)$ (ce qui établira que $X$ est normal, en vertu du critère de Serre \sref{IV.5.8.6}).
Raisonnons par l'absurde en supposant que l'ensemble $E$ des points $x\in X$ où $X$ ne vérifie pas $(\mathrm{S}_2)$ soit non vide, et soit $x\in E$ un point pour lequel $\dim(\sh{O}_{X,x})$ est la plus petite possible;
puisque $X$ vérifie $(\mathrm{S}_1)$, on a $\dim(\sh{O}_{X,x})\geq2$.
Dans $X_1=\Spec(\sh{O}_{X,x})$, l'ouvert $U=X_1-\{x\}$ vérifie $(\mathrm{S}_2)$;
montrons que $X_1$ vérifie $(\mathrm{S}_2)$, de sorte que l'on aura abouti à une contradiction.
Il suffit, en vertu de \sref{IV.5.10.5}, de montrer que toute section $f$ de $\sh{O}_{X_1}$ au-dessus de $U$ se prolonge en une section de $\sh{O}_{X_1}$ au-dessus de $X_1$.
Comme $X_1$ est réduit et $U$ dense dans $X_1$, $U$ est schématiquement dense dans $X_1$, et par suite $f$ est la restriction à $U$ d'une section méromorphe régulière $g\in M(X_1)$.
En outre, comme $\dim(\sh{O}_{X,x})\geq2$, on a $\operatorname{cyc}(\operatorname{div}(g))\geq0$ puisque $g|U=f$;
comme l'homomorphisme $\operatorname{cyc}$ est injectif, on a nécessairement $\operatorname{div}(g)\geq0$, donc \sref{IV.21.6.4} $g$ est une section de $\sh{O}_{X_1}$ au-dessus de $X_1$, ce qui établit notre assertion.

Pour montrer que \emph{a$'$)} implique \emph{c)}, remarquons alors que la condition \emph{a$'$)} implique que l'homomorphisme canonique $\operatorname{cyc}:\operatorname{Div}(X)\to\mathfrak{Z}^1(X)$ est surjectif;
puisque $X$ est normal, il suffit d'appliquer \sref{IV.21.6.9}.

Prouvons maintenant que \emph{a$''$)} entraîne \emph{c)}.
Il résulte de \sref{IV.21.6.5} que \emph{a$''$)} entraîne que tout cycle premier $1$-codimensionnel sur $X$ est l'image par $\operatorname{cyc}$ d'un diviseur positif;
la première partie du raisonnement fait ci-dessus prouve donc encore que $X$ vérifie $(\mathrm{R}_1)$ et $(\mathrm{S}_1)$.
Il reste encore à voir que $X$ est normal (la fin du raisonnement étant alors la même), c'est-à-dire que si $x\in X$ est tel que $\dim(\sh{O}_{X,x})\geq2$, on a $\operatorname{prof}(\sh{O}_{X,x})\geq2$.
Or, si $y$ est une générisation de $x$ telle que $\overline{\{y\}}$ soit de codimension $1$ dans $X$, le sous-préschéma réduit $Y$ de $X$ ayant $\overline{\{y\}}$ pour espace
\oldpage[IV-4]{280}
sous-jacent est intègre, donc régulièrement immergé dans $X$ par hypothèse.
Cela entraîne qu'il existe un élément régulier $t$ de $\sh{O}_{X,x}$ tel que l'idéal $t\sh{O}_{X,x}$ soit l'idéal premier définissant le préschéma $Y_1$ image réciproque de $Y$ dans $X_1=\Spec(\sh{O}_{X,x})$.
Comme $\sh{O}_{X,x}/t\sh{O}_{X,x}$ est intègre, cela prouve que $\operatorname{prof}(\sh{O}_{X,x})\geq2$ et termine la preuve de \sref{IV.21.7.4}.
\end{proof}

\begin{remark}[21.7.5]
\label{IV.21.7.5}
On ne peut dans \sref{IV.21.7.4} remplacer la condition \emph{a$'$)} par la condition plus faible que tout cycle premier $1$-codimensionnel de $X$ soit l'image par $\operatorname{cyc}$ d'un diviseur positif.
C'est ce que montre l'exemple où $X$ est le schéma affine défini dans \sref{IV.14.1.5} (« réunion d'un plan et d'une droite qui le coupe ») \emph{qui n'est pas normal};
avec les notations introduites dans \sref{IV.14.1.5}, les cycles premiers $1$-codimensionnels de $X$ sont ceux du plan $X_1$ et les points fermés de la droite $X_2$ distincts du point commun à $X_1$ et $X_2$;
si $t_1$, $t_2$, $t_3$ sont les images de $T_1$, $T_2$, $T_3$ dans $A/\mathfrak a$, on voit donc que les cycles premiers $1$-codimensionnels de $X$ sont définis par les idéaux premiers monogènes de générateurs $P(t_1,t_2)$ ($P$ irréductible dans $K[T_1,T_2]$) ou $t_3-a$, avec $a\neq0$ dans $K$;
ces éléments étant réguliers dans $A/\mathfrak a$, tout cycle $1$-codimensionnel est bien l'image canonique d'un diviseur positif.
\end{remark}

\begin{corollary}[21.7.6]
\label{IV.21.7.6}
Soit $X$ un préschéma localement noethérien, réduit en chacun de ses points isolés.
Les conditions suivantes sont équivalentes:
\begin{enumerate}
  \item[\emph{a)}] L'homomorphisme canonique $c:\sh{D}iv_X\to\mathscr{Z}_X^1$ est un isomorphisme de faisceaux de groupes ordonnés, et on a $\operatorname{Div}(X)=\operatorname{Div.princ}(X)$.
  \item[\emph{a$'$)}] L'homomorphisme canonique $c:\sh{D}iv_X\to\mathscr{Z}_X^1$ est injectif, et tout cycle premier $1$-codimensionnel est l'image par $\operatorname{cyc}$ d'un diviseur positif principal, autrement dit est de la forme $\operatorname{cyc}(\operatorname{div}(f))$, où $f$ est une section régulière de $\sh{O}_X$ au-dessus de $X$.
  \item[\emph{a$''$)}] Pour tout sous-préschéma fermé intègre $Y$ de $X$, de codimension $1$, il existe une section régulière $f$ de $\sh{O}_X$ au-dessus de $X$ telle que $Y$ soit défini par l'Idéal $f\sh{O}_X$.
  \item[\emph{b)}] $X$ est normal et tout cycle premier $1$-codimensionnel sur $X$ est principal.
  \item[\emph{c)}] $X$ est localement factoriel et $\operatorname{Pic}(X)=0$.
  \item[\emph{d)}] $X$ est normal, et pour tout ouvert $U$ de $X$, on a $\operatorname{Pic}(U)=0$.
\end{enumerate}
\end{corollary}

\begin{proof}
L'équivalence de \emph{a)}, \emph{a$'$)}, \emph{a$''$)}, \emph{b)} et \emph{c)} résulte aussitôt de \sref{IV.21.7.4} et de \sref{IV.21.6.11}.
En outre, il résulte aussitôt de \emph{a$''$)} que tout ouvert non vide $U$ de $X$ vérifie encore les mêmes conditions, autrement dit ces conditions impliquent \emph{d)}.
Reste à voir que \emph{d)} entraîne \emph{b)}.
Or, désignons par $U_\lambda$ les ouverts complémentaires des réunions finies d'ensembles de la forme $\overline{\{x\}}$, où $\dim(\sh{O}_{X,x})\geq2$;
il est clair que les $U_\lambda$ forment une famille filtrante décroissante et que pour tout $x\in\bigcap_\lambda U_\lambda$, on a $\dim(\sh{O}_{X,x})\leq1$, donc $\sh{O}_{X,x}$ est un anneau principal en vertu de l'hypothèse.
On peut donc appliquer à la famille $(U_\lambda)$ la proposition \sref{IV.21.6.12}, qui montre que $\mathfrak{Cl}(X)$ est isomorphe à $\varinjlim\operatorname{Pic}(U_\lambda)$, donc $\mathfrak{Cl}(X)=0$ en vertu de l'hypothèse, ce qui prouve \emph{b)} par définition.
\end{proof}

\begin{remark}[21.7.7]
\label{IV.21.7.7}
Lorsque $X=\Spec(A)$, où $A$ est un anneau noethérien intègre, les conditions équivalentes de \sref{IV.21.7.6} équivalent à dire que $A$ est un anneau \emph{factoriel}.
\end{remark}

\subsection{Diviseurs et normalisation}
\label{subsection:IV.21.8}

\begin{lemma}[21.8.1]
\label{IV.21.8.1}
Soit $f:X'\to X$ un morphisme entier de préschémas.
Pour tout $\sh{O}_{X'}$-Module localement libre $E'$ de rang constant $n$ et tout $x\in X$, il existe un voisinage ouvert $U$ de $x$ dans $X$ tel que, si l'on pose $U'=f^{-1}(U)$, $E'|U'$ soit isomorphe à $\sh{O}_{U'}^n$.
\end{lemma}

\begin{proof}
La question étant locale sur $X$, on peut se borner au cas où $X=\Spec(A)$ est affine, auquel cas on a $X'=\Spec(A')$, où $A'$ est une $A$-algèbre entière.
Alors $A'$ est limite inductive de ses sous-$A$-algèbres \emph{finies} $A'_\lambda$.
Posant $X'_\lambda=\Spec(A'_\lambda)$ et désignant par $p_\lambda:X'\to X'_\lambda$ le morphisme structural, il résulte de \sref{IV.8.5.2} et \sref{IV.8.5.5} qu'il existe un indice $\lambda$ et un $\sh{O}_{X'_\lambda}$-Module localement libre $E'_\lambda$ de rang $n$ tel que $E'=p_\lambda^*(E'_\lambda)$ et il suffira évidemment de prouver le lemme pour $X'_\lambda$ et $E'_\lambda$;
autrement dit, on peut se borner au cas où le morphisme $f$ est \emph{fini}.
Posons $B=\sh{O}_{X,x}$ et soit $B'$ l'anneau du schéma affine $X'_0=X'\times_X\Spec(\sh{O}_{X,x})$;
comme $B$ est un anneau local et que $B'$ est une $B$-algèbre finie, $B'$ est un anneau semi-local (Bourbaki, \emph{Alg.\ comm.}, chap.~V, \textsection~2, n\textsuperscript{o}~1, prop.~3);
on en conclut que le $\sh{O}_{X'_0}$-Module localement libre $E'_0=E'\otimes_{\sh{O}_{X'}}\sh{O}_{X'_0}$ est isomorphe à $\sh{O}_{X'_0}^n$ (Bourbaki, \emph{Alg.\ comm.}, chap.~II, \textsection~5, n\textsuperscript{o}~3, prop.~5).
Considérant $X'_0$ comme limite
\oldpage[IV-4]{281}
projective des préschémas induits par $X'$ sur les ouverts $f^{-1}(U)$, où $U$ parcourt l'ensemble filtrant des voisinages ouverts affines de $x$ dans $X$, suivant la méthode de \sref{IV.8.1.2}, \emph{a)}), et appliquant \sref{IV.8.5.2.5}, on obtient la conclusion cherchée.
\end{proof}

\begin{corollary}[21.8.2]
\label{IV.21.8.2}
Soit $f:X'\to X$ un morphisme entier de préschémas.
Alors:
\begin{enumerate}[label=(\roman*)]
  \item On a $\RR^1 f_*(\sh{O}^*_{X'})=0$ \sref[0]{0.12.2.1}.
  \item Le groupe $\operatorname{Pic}(X')$ est canoniquement isomorphe à $\HH^1(X,f_*(\sh{O}^*_{X'}))$.
\end{enumerate}
\end{corollary}

\begin{proof}
\begin{enumerate}[label=(\roman*)]
  \item $\RR^1 f_*(\sh{O}^*_{X'})$ est le faisceau associé au préfaisceau de groupes commutatifs $U\mapsto\HH^1(f^{-1}(U),\sh{O}^*_{X'})$ sur $X$ (\emph{loc.\ cit.});
la fibre de ce faisceau en un point $x$ peut donc \sref[0]{0.5.4.7} être identifiée au groupe commutatif $\varinjlim\operatorname{Pic}(f^{-1}(U))$, où $U$ parcourt l'ensemble des voisinages ouverts de $x$, les homomorphismes de transition $\operatorname{Pic}(f^{-1}(U'))\to\operatorname{Pic}(f^{-1}(U))$ pour deux ouverts $U\subset U'$ étant définis par \sref{IV.21.3.2.4}.
Or, pour tout $x\in X$, tout voisinage ouvert $U'$ de $x$ dans $X$, et tout élément $\zeta\in\operatorname{Pic}(f^{-1}(U'))$, il résulte de \sref{IV.21.8.1} qu'il y a un voisinage ouvert $U\subset U'$ de $x$ tel que l'image de $\zeta$ dans $\operatorname{Pic}(f^{-1}(U))$ soit nulle.
Donc la fibre de $\RR^1f_*(\sh{O}^*_{X'})$ au point $x$ est nulle.

  \item La suite spectrale de Leray pour le foncteur composé $\sh{F}'\mapsto\Gamma(X,f_*(\sh{F}'))$ de faisceaux de groupes commutatifs sur $X'$ (T, 2.4) donne la suite exacte de termes de bas degré
\[
  0\to\HH^1(X,f_*(\sh{O}^*_{X'}))\to\HH^1(X',\sh{O}^*_{X'})\to\HH^0(X,\RR^1f_*(\sh{O}^*_{X'}))
\]
et la conclusion résulte de (i) et de l'isomorphisme de $\operatorname{Pic}(X')$ et de $\HH^1(X',\sh{O}^*_{X'})$ \sref[0]{0.5.4.7}.
\end{enumerate}
\end{proof}

\begin{proposition}[21.8.3]
\label{IV.21.8.3}
Soit $f=(\psi,\theta):X'\to X$ un morphisme entier de préschémas;
supposons que, pour tout ouvert $U$ de $X$, l'homomorphisme $\Gamma(\theta):\Gamma(U,\sh{O}_X)\to\Gamma(U,f_*(\sh{O}_{X'}))$ transforme les éléments réguliers en éléments réguliers, ce qui permet de prolonger canoniquement l'homomorphisme $\theta:\sh{O}_X\to f_*(\sh{O}_{X'})$ en un homomorphisme de faisceaux d'anneaux $\theta':\sh{M}_X\to f_*(\sh{M}_{X'})$, d'où des homomorphismes de faisceaux de groupes multiplicatifs $\theta^*:\sh{O}^*_X\to f_*(\sh{O}^*_{X'})$ et $\theta'^*:\sh{M}^*_X\to f_*(\sh{M}^*_{X'})$, donnant par passage aux quotients un homomorphisme $\theta''^*:\sh{D}iv_X\to f_*(\sh{D}iv_{X'})$.
On a alors un diagramme commutatif
\[
\label{IV.21.8.3.1}
  \xymatrix{
    1 \ar[r] & \sh{O}^*_X \ar[r] \ar[d]^{\theta^*} & \sh{M}^*_X \ar[r] \ar[d]^{\theta'^*} & \sh{D}iv_X \ar[r] \ar[d]^{\theta''^*} & 0\\
    1 \ar[r] & f_*(\sh{O}^*_{X'}) \ar[r] & f_*(\sh{M}^*_{X'}) \ar[r] & f_*(\sh{D}iv_{X'}) \ar[r] & 0
  }
  \tag{21.8.3.1}
\]
dans lequel les deux lignes sont exactes.
\end{proposition}

\begin{proof}
Le seul point à vérifier est l'exactitude de la seconde ligne du diagramme, qui résulte de l'application de la suite exacte de cohomologie pour le foncteur $f_*$ à la suite exacte de faisceaux de groupes commutatifs sur $X'$
\[
  1\to\sh{O}^*_{X'}\to\sh{M}^*_{X'}\to\sh{D}iv_{X'}\to0
\]
et de \sref{IV.21.8.2}.
\end{proof}

\begin{corollary}[21.8.4]
\label{IV.21.8.4}
Si, en outre des hypothèses de \sref{IV.21.8.3}, l'homomorphisme $\theta'$ est un isomorphisme de faisceaux d'anneaux, alors $\theta^*:\sh{O}^*_X\to f_*(\sh{O}^*_{X'})$ est injectif, $\theta''^*:\sh{D}iv_X\to f_*(\sh{D}iv_{X'})$ est surjectif et $\operatorname{Ker}(\theta''^*)$ est isomorphe à $\operatorname{Coker}(\theta^*)$.
\end{corollary}

\begin{proof}
C'est une conséquence immédiate du diagramme du serpent (Bourbaki, \emph{Alg.\ comm.}, chap.~I, \textsection~2, n\textsuperscript{o}~4, prop.~2) appliqué au diagramme \sref{IV.21.8.3.1}.
\end{proof}

\oldpage[IV-4]{282}

\begin{proposition}[21.8.5]
\label{IV.21.8.5}
Soient $X$, $X'$ deux préschémas localement noethériens, $f:X'\to X$ un morphisme entier;
on suppose qu'il existe un ouvert schématiquement dense $U$ dans $X$ tel que $f^{-1}(U)$ soit schématiquement dense dans $X'$ (cf.\ \sref{IV.20.3.5}), et que le morphisme $f^{-1}(U)\to U$, restriction de $f$, soit un isomorphisme.
Alors:
\begin{enumerate}[label=(\roman*)]
  \item L'homomorphisme $\theta':\sh{M}_X\to f_*(\sh{M}_{X'})$ est défini et bijectif, l'homomorphisme
\[
  \theta^*:\sh{O}^*_X\to f_*(\sh{O}^*_{X'})
\]
est injectif, l'homomorphisme $\theta''^*:\sh{D}iv_X\to f_*(\sh{D}iv_{X'})$ est surjectif, $\operatorname{Ker}(\theta''^*)$ est isomorphe à $\sh{N}=\operatorname{Coker}(\theta^*)$ et le support du faisceau de groupes multiplicatifs $\sh{N}$ est contenu dans $S=X-U$.
  \item Supposons de plus l'ensemble $S$ discret et posons pour abréger $\sh{O}'_X=f_*(\sh{O}_{X'})$.
Alors on a un diagramme commutatif
\[
\label{IV.21.8.5.1}
  \xymatrix{
    1 \ar[r] & \prod_{s\in S}\sh{O}'^{*}_{X,s}/\sh{O}^*_{X,s} \ar[r]^-{j} \ar[d] & \operatorname{Div}(X) \ar[r] \ar[d]^{k_X} & \operatorname{Div}(X') \ar[r] \ar[d]^{k_{X'}} & 0\\
    1 \ar[r] & \bigl(\prod_{s\in S}\sh{O}'^{*}_{X,s}/\sh{O}^*_{X,s}\bigr)/\operatorname{Im}\Gamma(X',\sh{O}^*_{X'}) \ar[r]^-{j'} & \operatorname{Pic}(X) \ar[r] & \operatorname{Pic}(X') \ar[r] & 0
  }
  \tag{21.8.5.1}
\]
où les lignes sont exactes et où la flèche verticale de gauche est l'homomorphisme canonique.
\end{enumerate}
\end{proposition}

\begin{proof}
\begin{enumerate}[label=(\roman*)]
  \item L'hypothèse entraîne que l'on a un isomorphisme canonique $\sh{M}''_X\simto f_*(\sh{M}''_{X'})$ pour les faisceaux de germes de pseudo-fonctions \sref{IV.20.2.2}, d'où l'assertion relative à $\theta'$ vu l'existence des isomorphismes canoniques $\sh{M}_X\simto\sh{M}''_X$ et $\sh{M}_{X'}\simto\sh{M}''_{X'}$ \sref{IV.20.2.11}.
Le reste de l'assertion (i) découle de \sref{IV.21.8.4}, sauf ce qui concerne le support de $\sh{N}$, qui résulte directement de l'hypothèse sur $U$.

  \item Appliquant la suite exacte de cohomologie aux deux suites exactes de faisceaux de groupes commutatifs
\[
\label{IV.21.8.5.2}
\begin{gathered}
  1\to\sh{N}\to\sh{D}iv_X\to f_*(\sh{D}iv_{X'})\to0,\\
  1\to\sh{O}^*_X\to f_*(\sh{O}^*_{X'})\to\sh{N}\to1,
\end{gathered}
  \tag{21.8.5.2}
\]
il vient respectivement les suites exactes
\[
\label{IV.21.8.5.3}
\begin{gathered}
  1\to\Gamma(X,\sh{N})\xrightarrow{j}\operatorname{Div}(X)\to\operatorname{Div}(X')\to\HH^1(X,\sh{N}),\\
  \Gamma(X',\sh{O}^*_{X'})\to\Gamma(X,\sh{N})\xrightarrow{\partial}\operatorname{Pic}(X)\to\operatorname{Pic}(X')\to\HH^1(X,\sh{N}),
\end{gathered}
  \tag{21.8.5.3}
\]
\oldpage[IV-4]{283}
compte tenu de l'isomorphisme canonique $\operatorname{Pic}(X')\simto\HH^1(X,f_*(\sh{O}^*_{X'}))$ \sref{IV.21.8.2}.
Or, comme le support de $\sh{N}$ est contenu dans l'ensemble discret $S$, fermé dans $X$, on a
\[
  \HH^1(X,\sh{N})=\HH^1(S,\sh{N}|S)\quad\text{(G, II, 4.9.2)},\qquad
  \HH^1(S,\sh{N}|S)=\prod_{s\in S}\HH^1(\{s\},\sh{N}_s)=0
\]
par définition de la cohomologie (G, II, 4.4).
De même on a $\Gamma(X,\sh{N})=\prod_{s\in S}\sh{N}_s$ et $\sh{N}_s=\sh{O}'^{*}_{X,s}/\sh{O}^*_{X,s}$ en vertu de la seconde suite exacte \sref{IV.21.8.5.2}.
Reste à préciser les injections $j$ et $j'$.
Or, on peut décrire une section $t$ de $\sh{N}$ au-dessus de $X$ en prenant un recouvrement de $X$ formé de $U$ et d'ouverts $U_i$ tels que $U_i\cap S$ soit réduit à un seul point $s_i$, et que $U_i\cap U_j\subset U$ pour $i\neq j$, puis en considérant pour chaque $i$ une section $t_i$ de $\sh{N}$ au-dessus de $U_i$, nécessairement telle que $(t_i)_{s_i}\in\sh{O}'^{*}_{X,s_i}/\sh{O}^*_{X,s_i}$ et $(t_i)_x=0$ aux points $x\in U_i-\{s_i\}$.
Le germe $(t_i)_{s_i}$ provient d'une section $u_i$ de $\sh{O}^*_{X'}$ au-dessus de $f^{-1}(U_i)$, qu'on peut aussi considérer comme une section de $\sh{M}^*_{X'}$ au-dessus de $f^{-1}(U_i)$, donc une section de $\sh{M}^*_X$ au-dessus de $U_i$ (en vertu de l'hypothèse);
il correspond canoniquement à cette section une section $d_i\in\Gamma(U_i,\sh{D}iv_X)$ et comme dans $f^{-1}(U\cap U_i)$ la restriction de $u_i$ s'identifie à une section de $\sh{O}^*_X$ au-dessus de $U\cap U_i$ en vertu de l'hypothèse, les restrictions des $d_i$ à $U\cap U_i$ sont toutes nulles, donc les $d_i$ sont les restrictions d'un même diviseur de $\operatorname{Div}(X)$, qui est l'image de la section $t$ par $j$.
De même, l'image de $t$ par $\partial:\Gamma(X,\sh{N})\to\HH^1(X,\sh{O}^*_X)$ provient du cocycle égal à la restriction de $u_i$ dans $U\cap U_i$ pour chaque $i$, à $(u_i|(U_i\cap U_j))(u_j|(U_i\cap U_j))^{-1}$ dans $U_i\cap U_j$, d'où l'expression de $j'(t)$ et la commutativité du diagramme \sref{IV.21.8.5.1}.
\end{enumerate}
\end{proof}

\begin{remark}[21.8.6]
\label{IV.21.8.6}
\begin{enumerate}[label=(\roman*)]
  \item Les conditions d'application de \sref{IV.21.8.5}, (i) sont en particulier remplies lorsque $X$ est un préschéma localement noethérien réduit n'ayant qu'un nombre fini de composantes irréductibles, $X'$ son normalisé et que $X'$ est aussi localement noethérien (\sref[II]{II.6.3.8});
le $\sh{O}_X$-Module $\sh{D}iv_X$ est alors une \emph{extension} de $f_*(\sh{D}iv_{X'})$ par le conoyau $\sh{N}$ de $\sh{O}^{*}_X\to f_*(\sh{O}^*_{X'})$, que l'on peut considérer comme connu.
Si de plus $X$ est une courbe algébrique (réduite) sur un corps $k$, on est dans les conditions d'application de \sref{IV.21.8.5}, (ii).

  \item Lorsque les conditions d'application de \sref{IV.21.8.5}, (ii) sont remplies et en outre que l'homomorphisme canonique global $\Gamma(X,\sh{O}_X)\to\Gamma(X',\sh{O}_{X'})$ est \emph{bijectif}, on voit que les noyaux des homomorphismes surjectifs $\operatorname{Div}(X)\to\operatorname{Div}(X')$ et $\operatorname{Pic}(X)\to\operatorname{Pic}(X')$ sont \emph{isomorphes}.

  \item Lorsque les conditions d'application de \sref{IV.21.8.5}, (ii) sont remplies, on voit que l'homomorphisme $\operatorname{Div}(X)\to\operatorname{Div}(X')$ ne peut être injectif (auquel cas il est bijectif) que si $\sh{O}'^{*}_{X,s}=\sh{O}^*_{X,s}$ pour tout $s\in S$.
Pour un $s\in S$ tel que l'anneau $\sh{O}'_{X,s}$ soit \emph{local} (ce qui, compte tenu du fait que $X'=\Spec(\sh{O}'_X)$ \sref[II]{II.1.3}, signifie qu'il n'existe qu'un seul point $s'\in X'$ au-dessus de $s$) cela entraîne nécessairement que les corps résiduels $\boldsymbol{k}(s)$ et $\boldsymbol{k}(s')$ sont égaux et que $1+\mathfrak{m}_s=1+\mathfrak{m}_{s'}$, donc finalement équivaut à la relation $\sh{O}'_{X,s}=\sh{O}_{X,s}$ ou encore (compte tenu de l'hypothèse) au fait qu'il y a un voisinage ouvert $V$ de $s$ dans $X$ tel que le morphisme $f^{-1}(V)\to V$ soit un isomorphisme.
Dans le cas général, l'anneau $\sh{O}'_{X,s}$ est semi-local (c'est évident lorsque le morphisme $f$ est fini et dans le cas général cela résulte de \sref[0]{0.23.2.6}), l'homomorphisme canonique $\sh{O}^{*}_{X,s}\to\sh{O}'^{*}_{X,s}$ définit, par passage aux quotients, un homomorphisme du groupe multiplicatif $\boldsymbol{k}(s)^*$ dans le produit des groupes multiplicatifs $\boldsymbol{k}(s'_j)^*$, les $s'_j$ étant les points (en
\oldpage[IV-4]{284}
nombre $>1$) de $X'$ au-dessus de $s$.
Il est immédiat que cet homomorphisme ne peut être bijectif que si $\boldsymbol{k}(s)$ et tous les $\boldsymbol{k}(s'_j)$ sont égaux au corps $\bb{F}_2$ à $2$ éléments;
de plus, si cette condition est vérifiée, il faut en outre que le groupe multiplicatif $1+\mathfrak{m}_s$ ait pour image $1+\mathfrak{r}$, où $\mathfrak{r}$ est le radical de $\sh{O}'_{X,s}$, ou encore que $\mathfrak{m}_s=\mathfrak{r}$, ce qui entraîne que $\sh{O}'_{X,s}\otimes_{\sh{O}_{X,s}}\boldsymbol{k}(s)$ est composé direct de corps isomorphes à $\boldsymbol{k}(s)$ (ce qui, lorsque le morphisme $f$ est fini, entraîne qu'il est non ramifié aux points $s'_j$ \sref{IV.17.4.1}).
Si par exemple \emph{aucun} corps résiduel de $X$ n'est isomorphe à $\bb{F}_2$, l'homomorphisme canonique $\operatorname{Div}(X)\to\operatorname{Div}(X')$ ne peut être bijectif que si $f$ est un \emph{isomorphisme}.
Dans le cas où l'homomorphisme canonique $\Gamma(X,\sh{O}_X)\to\Gamma(X',\sh{O}_{X'})$ est bijectif, on conclut de ce qui précède et de (ii) que dans les considérations précédentes, on peut remplacer l'homomorphisme $\operatorname{Div}(X)\to\operatorname{Div}(X')$ par l'homomorphisme $\operatorname{Pic}(X)\to\operatorname{Pic}(X')$.
\end{enumerate}
\end{remark}

\subsection{Diviseurs sur les préschémas de dimension 1}
\label{subsection:IV.21.9}
\label{IV.21.9}

\begin{env}[21.9.1]
\label{IV.21.9.1}
Soient $X$ un espace topologique, $x$ un point de $X$, $i_x:\{x\}\to X$ l'injection canonique.
Si $A(x)$ est un groupe commutatif, on peut le considérer comme un faisceau de groupes commutatifs sur l'espace $\{x\}$ réduit à un seul point, et prendre son image directe $(i_x)_*(A(x))$ qui est un faisceau de groupes commutatifs sur $X$ \sref[0]{0.3.4.1};
il résulte aussitôt des définitions que pour tout ouvert $U$ de $X$, $\Gamma(U,(i_x)_*(A(x)))=A(x)$ si $x\in U$, et est réduit à $0$ dans le cas contraire;
donc, pour $y\in\overline{\{x\}}$ on a $((i_x)_*(A(x)))_y=A(x)$, et pour $y\notin\overline{\{x\}}$, $((i_x)_*(A(x)))_y=0$.

Si maintenant $\sh{F}$ est un faisceau de groupes commutatifs sur $X$ et $Y$ une partie de $X$, pour tout ouvert $U$ de $X$, on a un homomorphisme canonique
\[
\label{IV.21.9.1.1}
  \Gamma(U,\sh{F})\to\prod_{x\in U\cap Y}\sh{F}_x
  =\prod_{x\in Y}\Gamma(U,(i_x)_*(\sh{F}_x))
  \tag{21.9.1.1}
\]
et comme ces homomorphismes commutent aux opérateurs de restriction, ils définissent un homomorphisme canonique de \emph{faisceaux}
\[
\label{IV.21.9.1.2}
  j_Y:\sh{F}\to\prod_{x\in Y}(i_x)_*(\sh{F}_x).
  \tag{21.9.1.2}
\]
\end{env}

\begin{lemma}[21.9.2]
\label{IV.21.9.2}
Soient $X$ un espace topologique localement noethérien, $X_0$ l'ensemble de ses points fermés, $\sh{F}$ un faisceau de groupes commutatifs sur $X$.
Les conditions suivantes sont équivalentes:
\begin{enumerate}[label=\emph{\alph*)}]
  \item L'homomorphisme canonique $j_{X_0}:\sh{F}\to\prod_{x\in X_0}(i_x)_*(\sh{F}_x)$ est injectif et a pour image $\bigoplus_{x\in X_0}(i_x)_*(\sh{F}_x)$.
  \item[\rm a$'$)] Il existe une famille de groupes commutatifs $(A(x))_{x\in X_0}$ telle que $\sh{F}$ soit isomorphe à $\bigoplus_{x\in X_0}(i_x)_*(A(x))$.
  \item Toute section de $\sh{F}$ au-dessus d'un ouvert $U$ de $X$ a un support discret contenu dans $X_0$.
\end{enumerate}
Lorsqu'il en est ainsi, pour tout $x\in X_0$, le groupe $A(x)$ est déterminé à isomorphisme unique près et est isomorphe à $\sh{F}_x$.
En outre, le faisceau $\sh{F}$ est alors flasque.
\end{lemma}

\begin{proof}
\oldpage[IV-4]{285}
Comme les points de $X_0$ sont fermés dans $X$, on a $((i_x)_*(A(x)))_y=0$ pour tout $y\neq x\in X_0$;
cette remarque prouve l'assertion d'unicité relative aux groupes $A(x)$, et il est clair par ailleurs que \emph{a)} implique \emph{a$'$)}.
Pour voir que \emph{a$'$)} implique \emph{b)}, on peut se borner au cas où $U$ est noethérien;
alors (G, II, 3.10) on a
\[
  \Gamma(U,\bigoplus_{x\in X_0}(i_x)_*(A(x)))
  =\bigoplus_{x\in X_0}\Gamma(U,(i_x)_*(A(x))),
\]
donc toute section $s$ de $\sh{F}$ au-dessus de $U$ est somme directe d'un nombre \emph{fini} de sections $s_k\in\Gamma(U,(i_{x_k})_*(A(x_k)))$ ($1\leq k\leq n$) avec $x_k\in X_0$.
Mais puisque $x_k$ est fermé dans $X$, le support de $s_k$ est réduit à $\{x_k\}$, donc le support de $s$ est contenu dans l'ensemble fini des $x_k$, qui est évidemment discret puisque les points $x_k$ sont fermés dans $X$.
Montrons enfin que \emph{b)} entraîne \emph{a)};
pour tout ouvert noethérien $U$, tout support d'une section de $\sh{F}$ au-dessus de $U$ étant discret et quasi-compact, est \emph{fini}.
On déduit donc de l'hypothèse \emph{b)} que pour tout ouvert noethérien $U$, l'image de l'homomorphisme \sref{IV.21.9.1.1} (pour $Y=X_0$) est contenue dans $\bigoplus_{x\in X_0}\Gamma(U,(i_x)_*(\sh{F}_x))$ et que cet homomorphisme est injectif, ce qui établit \emph{a)}.

Enfin, pour montrer que $\sh{F}$ est flasque, considérons une section $s$ de $\sh{F}$ au-dessus d'un ouvert $U$ de $X$;
comme le support de $s$ est un sous-espace discret et fermé dans $X$, on prolonge $s$ en une section $s'$ de $\sh{F}$ au-dessus de $X$ en posant $s'_x=0$ dans $X-U$.
\end{proof}

\begin{remark}[21.9.3]
\label{IV.21.9.3}
\begin{enumerate}[label=(\roman*)]
  \item Dans la condition \emph{b)} de \sref{IV.21.9.2}, il ne suffit pas de supposer que le support de toute section de $\sh{F}$ au-dessus d'un ouvert quelconque de $X$ soit \emph{discret};
c'est ce que montre l'exemple où on prend pour $X$ un spectre d'anneau de valuation discrète, avec un point fermé $b$ et un point générique $a$, et pour $\sh{F}$ le faisceau de groupes commutatifs tel que $\sh{F}_b=0$ et $\sh{F}_a=\bb{Z}$.

  \item Les conditions de \sref{IV.21.9.2} étant supposées vérifiées, soit $E$ une partie discrète de $X_0$, et pour tout $x\in E$, soit $a(x)$ un élément de $A(x)$.
Il existe alors une section $t$ de $\bigoplus_{x\in X_0}(i_x)_*(A(x))$ et une seule au-dessus de $X$ telle que $t_x=a(x)$ pour tout $x\in E$ et que le support de $t$ soit contenu dans $E$.
En effet, pour tout ouvert noethérien $U$, $E\cap U$ est fini, et il suffit de prendre pour $t$ la section de $\bigoplus_{x\in E}(i_x)_*(A(x))$ dont la restriction à chaque ouvert noethérien $U$ est la somme des $a(x)$ pour $x\in E\cap U$.
\end{enumerate}
\end{remark}

\begin{proposition}[21.9.4]
\label{IV.21.9.4}
Soient $X$ un préschéma localement noethérien de dimension $\leq1$, $X^{(1)}$ l'ensemble des points $x\in X$ tels que $\dim(\sh{O}_{X,x})=1$.
On a alors un isomorphisme canonique
\[
\label{IV.21.9.4.1}
  \sh{D}iv_X\simto\bigoplus_{x\in X^{(1)}}(i_x)_*(\operatorname{Div}(\sh{O}_{X,x}))
  \tag{21.9.4.1}
\]
et $\sh{D}iv_X$ est flasque.
\end{proposition}

\begin{proof}
Compte tenu de l'isomorphisme \sref{IV.21.4.6.1}, l'homomorphisme \sref{IV.21.9.4.1} est défini par \sref{IV.21.9.1.2}:
prouvons que c'est une bijection.
Comme $\dim(X)\leq1$, les points de $X^{(1)}$ sont les \emph{points fermés non isolés} de $X$.
On est ramené à prouver que:
1\textsuperscript{o} $\sh{D}iv_X$ vérifie la condition \emph{b)} de \sref{IV.21.9.2};
2\textsuperscript{o} pour tout point isolé $x\in X$, on a $(\sh{D}iv_X)_x=0$.
Le second point résulte de ce qu'alors $\sh{O}_{X,x}$ est un anneau artinien et de ce que tout élément régulier de $\sh{O}_{X,x}$ est donc inversible.
\oldpage[IV-4]{286}
Pour le 1\textsuperscript{o}, il suffit de noter que pour tout ouvert $U$ de $X$ et tout diviseur $D\in\operatorname{Div}(U)$, les points maximaux du support $S$ de $D$ sont tels que $\operatorname{prof}(\sh{O}_{X,x})=1$ \sref{IV.21.1.9}, donc \emph{a fortiori} appartiennent à $X^{(1)}\cap U$;
puisque $\dim(X)\leq1$, l'ensemble de ces points est égal à $S$, et $S$ est donc discret, l'ensemble des composantes irréductibles de $S$ étant localement fini.
\end{proof}

\begin{corollary}[21.9.5]
\label{IV.21.9.5}
Soit $X$ un préschéma localement noethérien de dimension $\leq1$.
Pour toute partie discrète $E\subset X^{(1)}$ et toute famille $(D(x))_{x\in E}$ avec $D(x)\in\operatorname{Div}(\sh{O}_{X,x})$, il existe un diviseur $D$ sur $X$ et un seul tel que le support de $D$ soit contenu dans $E$ et que $D_x=D(x)$ pour tout $x\in E$.
\end{corollary}

\begin{proof}
Cela résulte de \sref{IV.21.9.4} et de \sref{IV.21.9.3}, (ii).
\end{proof}

\begin{corollary}[21.9.6]
\label{IV.21.9.6}
Soit $X$ un préschéma localement noethérien de dimension $\leq1$;
tout diviseur $D$ sur $X$ est de la forme $D'-D''$, où $D'$, $D''$ sont deux diviseurs $\geq0$ de supports contenus dans celui de $D$.
\end{corollary}

\begin{proof}
En vertu de \sref{IV.21.9.5}, on est aussitôt ramené au cas où $X=\Spec(A)$, où $A$ est un anneau local noethérien;
il suffit alors d'observer que $M(X)$ est l'anneau total des fractions de $A$, et qu'une section de $\sh{M}^*_X$ au-dessus de $X$ s'écrit par définition comme quotient $b/a$ de deux éléments réguliers de $A$.
\end{proof}

\begin{corollary}[21.9.7]
\label{IV.21.9.7}
Soient $X$ un préschéma localement noethérien de dimension $\leq1$, sans point isolé, et $U$ un ouvert dense dans $X$.
Il existe alors un diviseur $D\geq0$ sur $X$, de support contenu dans $U$ et rencontrant chacune des composantes irréductibles de $X$.
\end{corollary}

\begin{proof}
On peut supposer que $U$ est réunion d'ensembles ouverts disjoints non vides $U_\alpha$, dont chacun est contenu dans une seule composante irréductible de $X$;
il suffit alors de prendre dans chaque $U_\alpha$ un point $x_\alpha$ fermé dans $X$ (il en existe puisque l'unique point générique de $U_\alpha$ ne peut être isolé par hypothèse), et d'appliquer \sref{IV.21.9.5} à l'ensemble discret des $x_\alpha$.
\end{proof}

L'intérêt de ce corollaire est qu'il permettra de prouver qu'une courbe algébrique séparée $X$ sur un corps $k$ est quasi-projective, le diviseur $D\geq0$ défini dans \sref{IV.21.9.7} étant alors nécessairement \emph{ample} en vertu du théorème de Riemann--Roch pour les courbes (chap.~V).

\begin{env}[21.9.8]
\label{IV.21.9.8}
Pour les préschémas localement noethériens $X$ de dimension $\leq1$ la proposition \sref{IV.21.9.4} ramène la détermination de $\sh{D}iv_X$ au cas où $X=\Spec(A)$, $A$ étant un anneau local noethérien de dimension $1$.

Lorsque $A$ est un anneau local régulier de dimension $1$ (autrement dit, un anneau de valuation discrète), le groupe $\operatorname{Div}(A)$ s'identifie canoniquement à $\bb{Z}$ \sref{IV.21.6.8};
par suite, en vertu de \sref{IV.21.9.2}:
\end{env}

\begin{proposition}[21.9.9]
\label{IV.21.9.9}
Soit $X$ un préschéma localement noethérien régulier de dimension $\leq1$.
Alors le faisceau $\sh{D}iv_X$ est canoniquement isomorphe à $\bigoplus_{x\in X^{(1)}}(i_x)_*(\bb{Z}(x))$, où $\bb{Z}(x)$ est le groupe additif $\bb{Z}$ considéré comme faisceau de groupes sur l'espace $\{x\}$.
\end{proposition}

\begin{env}[21.9.10]
\label{IV.21.9.10}
Supposons seulement que l'anneau local noethérien $A$ de dimension $1$ soit \emph{réduit};
alors, si $A'$ est le \emph{normalisé} de $A$ (fermeture intégrale de $A$ dans son anneau total de fractions), $A'$ est noethérien en vertu du théorème de Krull--Akizuki (Bourbaki, \emph{Alg.\ comm.}, chap.~VII, \textsection~2, n\textsuperscript{o}~5, prop.~5), et on a vu dans \sref{IV.21.8.6} comment
\oldpage[IV-4]{287}
$\operatorname{Div}(A)$ peut s'obtenir comme \emph{extension} de $\operatorname{Div}(A')$ et ce dernier groupe est de la forme $\bb{Z}^r$.
\end{env}

\begin{proposition}[21.9.11]
\label{IV.21.9.11}
Soient $X$ un préschéma noethérien, $X_0$ un sous-préschéma fermé de $X$ ayant les propriétés suivantes:
\begin{enumerate}
  \item[\rm 1\textsuperscript{o}] $\dim(X_0)\leq1$.
  \item[\rm 2\textsuperscript{o}] Pour toute partie localement fermée $Y$ de $X$ telle que $Y\cap X_0$ soit discret, il existe une partie $Y'$ de $Y$, fermée dans $X$ et ouverte dans $Y$, contenant $Y\cap X_0$.
\end{enumerate}
Dans ces conditions:
\begin{enumerate}[label=(\roman*)]
  \item Soit $Z_0$ la réunion des ensembles $\overline{\{x\}}\cap X_0$ lorsque $x$ parcourt la partie de $\operatorname{Ass}(\sh{O}_X)$ formée des points tels que $\overline{\{x\}}\cap X_0$ soit fini.
Alors, pour tout diviseur $D_0$ sur $X_0$, dont le support ne rencontre pas $Z_0$, il existe un diviseur $D$ sur $X$ dont l'image réciproque par l'injection canonique $X_0\to X$ existe \sref{IV.21.4} et est égale à $D_0$;
si de plus $D_0\geq0$, on peut supposer $D\geq0$.

  \item Supposons en outre qu'il existe dans $X_0$ un ouvert affine $U_0$ contenant $\operatorname{Ass}(\sh{O}_{X_0})\cup Z_0$ (condition automatiquement vérifiée lorsqu'il existe un $\sh{O}_{X_0}$-Module ample \sref[II]{II.4.5.4}).
Alors l'homomorphisme canonique \sref{IV.21.3.2.4} $\operatorname{Pic}(X)\to\operatorname{Pic}(X_0)$ est surjectif.
\end{enumerate}
\end{proposition}

\begin{proof}
\begin{enumerate}[label=(\roman*)]
  \item Compte tenu de \sref{IV.21.9.6}, on peut se borner à prouver l'assertion relative aux diviseurs $D_0\geq0$;
en vertu de \sref{IV.21.9.4}, le support $T$ de $D_0$ est un ensemble discret et fermé dans $X_0$.
On peut supposer $D_0\neq0$, c'est-à-dire $T\neq\emp$, sans quoi il n'y a rien à démontrer.
Pour tout $x\in T$, il existe un élément $s_x\in\sh{O}_{X_0,x}$ qui est non diviseur de zéro dans cet anneau, appartient à son idéal maximal, et dont l'image dans $(\sh{D}iv_{X_0})_x$ est $(D_0)_x$.
Il existe un voisinage ouvert affine $U^{(x)}$ de $x$ dans $X$ et une section $g^{(x)}$ de $\sh{O}_X$ au-dessus de $U^{(x)}$ tels que $s_x$ soit l'image dans $\sh{O}_{X_0,x}$ du germe $(g^{(x)})_x\in\sh{O}_{X,x}$;
nous allons voir qu'en prenant $U^{(x)}$ assez petit, on peut faire en sorte que $g^{(x)}$ ne soit pas diviseur de zéro dans $\Gamma(U^{(x)},\sh{O}_X)$, autrement dit \sref{IV.3.1.9}, en désignant par $V^{(x)}$ la partie fermée de $U^{(x)}$ formée des $y$ tels que $g^{(x)}(y)=0$, que l'on ait $V^{(x)}\cap\operatorname{Ass}(\sh{O}_X)=\emp$.

Or, si $z\in\operatorname{Ass}(\sh{O}_X)$, l'hypothèse $x\notin Z_0$ entraîne que l'on a, ou bien $x\notin\overline{\{z\}}$, ou bien que $x$ n'est pas isolé dans $\overline{\{z\}}\cap X_0$.
En remplaçant $U^{(x)}$ par un ouvert plus petit, on peut supposer que c'est le second cas qui se produit pour un $z\in V^{(x)}\cap\operatorname{Ass}(\sh{O}_X)$.
$V^{(x)}$ contiendrait donc une composante irréductible de dimension $1$ de $X_0\cap U^{(x)}$, contenant $x$;
mais cela signifierait que l'image $\overline{g}^{(x)}$ de $g^{(x)}$ dans $\Gamma(U^{(x)}\cap X_0,\sh{O}_{X_0})$ serait dans le nilradical de cet anneau, et par suite $s_x$ appartiendrait au nilradical de $\sh{O}_{X_0,x}$, ce qui est absurde.
On a donc $x\notin\overline{\{z\}}$ pour tout point $z\in V^{(x)}\cap\operatorname{Ass}(\sh{O}_X)$, et comme cet ensemble est fini, on peut, en remplaçant $U^{(x)}$ par un voisinage ouvert plus petit de $x$, supposer que $V^{(x)}\cap\operatorname{Ass}(\sh{O}_X)=\emp$.

En vertu du choix de $U^{(x)}$, on peut définir un diviseur $D'^{(x)}$ sur $U^{(x)}$ par $D'^{(x)}=\operatorname{div}(g^{(x)})$;
d'ailleurs on a vu ci-dessus que $x$ est nécessairement isolé dans $V^{(x)}\cap X_0$, donc en remplaçant encore $U^{(x)}$ par un voisinage ouvert plus petit de $x$, on peut supposer que $V^{(x)}\cap X_0$ est réduit au point $x$.
Mais il existe alors, en vertu de la condition 2\textsuperscript{o}, une partie $W^{(x)}$ de $V^{(x)}$, fermée dans $X$ et ouverte dans $V^{(x)}$, telle que $W^{(x)}\cap X_0=\{x\}$.
Remplaçant encore $U^{(x)}$ par un voisinage ouvert plus petit de $x$, on peut donc supposer
\oldpage[IV-4]{288}
que $W^{(x)}=V^{(x)}$, autrement dit que $V^{(x)}$ est fermé dans $X$.
On peut alors définir un diviseur $D^{(x)}$ sur $X$ par la condition que $D^{(x)}|U^{(x)}=D'^{(x)}$, et $D^{(x)}|(X-V^{(x)})=0$, ce qui a un sens puisque dans $U^{(x)}\cap(X-V^{(x)})$ la restriction de $g^{(x)}$ est une section inversible, donc les restrictions à cet ouvert des deux diviseurs considérés sur $U^{(x)}$ et $X-V^{(x)}$ sont égales.
Il suffit alors, pour répondre à la question, de prendre $D=\sum_{x\in T}D^{(x)}$, ce qui a un sens puisque $T$ est fini.

  \item Compte tenu du diagramme commutatif \sref{IV.21.4.2.1}, la conclusion résultera de (i) si l'on prouve que tout $\sh{O}_{X_0}$-Module inversible $\sh{L}_0$ est isomorphe à un $\sh{O}_{X_0}$-Module de la forme $\sh{O}_{X_0}(D_0)$ \sref{IV.21.2.8}, où $D_0$ est un diviseur sur $X_0$ dont le support ne rencontre pas $Z_0$.
Or, comme $U_0$ est schématiquement dense dans $X_0$ \sref{IV.20.2.13}, (iv), il suffit pour cela de prouver qu'il existe une section $t\in\Gamma(U_0,\sh{L}_0)$ telle que $t(z_0)\neq0$ aux points de $\operatorname{Ass}(\sh{O}_{X_0})\cup Z_0$ \sref{IV.3.1.9}.
On peut donc se borner au cas où $X_0=\Spec(A_0)$ est affine;
mais alors $\sh{L}_0$ est ample \sref[II]{II.5.1.4} et comme l'ensemble $\operatorname{Ass}(\sh{O}_{X_0})\cup Z_0$ est fini, la conclusion résulte de \sref[II]{II.4.5.4}.
\end{enumerate}
\end{proof}

\begin{corollary}[21.9.12]
\label{IV.21.9.12}
Soient $A$ un anneau local hensélien \sref{IV.18.5.8}, $S=\Spec(A)$, $s_0$ le point fermé de $S$, $f:X\to S$ un morphisme séparé de présentation finie, et supposons que l'ensemble $X_0=f^{-1}(s_0)$ soit de dimension $\leq1$.
Alors, pour tout sous-schéma fermé $X'_0$ de $X$, ayant $X_0$ pour ensemble sous-jacent et de présentation finie sur $S$, l'homomorphisme canonique \sref{IV.21.3.2.4} $\operatorname{Pic}(X)\to\operatorname{Pic}(X'_0)$ est surjectif.
\end{corollary}

\begin{proof}
Montrons d'abord qu'il suffit de prouver le corollaire lorsque l'anneau local hensélien $A$ est de plus noethérien.
On sait en effet \sref{IV.18.6.15} que l'on peut écrire $A=\varinjlim A_\lambda$, où les $A_\lambda$ sont des anneaux locaux noethériens et henséliens, les homomorphismes $A_\lambda\to A$ étant locaux;
il existe en outre un indice $\alpha$ et un morphisme séparé de présentation finie $f_\alpha:X_\alpha\to S_\alpha=\Spec(A_\alpha)$ tels que $X$ et $f$ soient déduits de $X_\alpha$ et $f_\alpha$ par le changement de base $S\to S_\alpha$ \sref{IV.8.10.5}, (v);
avec les notations habituelles \sref{IV.8.8.1}, les morphismes $f_\lambda:X_\lambda\to S_\lambda$ seront séparés pour $\lambda\geq\alpha$ et on aura $X=X_\lambda\times_{S_\lambda}S$.
En outre, on peut supposer que, si $s_{0\alpha}$ est l'unique point fermé de $S_\alpha$, il y a un sous-schéma fermé $X'_{0\alpha}$ de $X_\alpha$, ayant même ensemble sous-jacent que $f_\alpha^{-1}(s_{0\alpha})$, tel que $X'_0=X'_{0\alpha}\times_{S_\alpha}S$ \sref{IV.8.6.3};
pour $\lambda\geq\alpha$, posons $X'_{0\lambda}=X'_{0\alpha}\times_{S_\alpha}S_\lambda$;
on a $\dim(X'_{0\lambda})=\dim(X'_0)\leq1$ par transitivité des fibres et \sref{IV.4.1.4}.
Il s'agit de voir que si l'on a prouvé que l'homomorphisme $\operatorname{Pic}(X_\lambda)\to\operatorname{Pic}(X'_{0\lambda})$ est surjectif pour tout $\lambda\geq\alpha$, il en est de même de $\operatorname{Pic}(X)\to\operatorname{Pic}(X'_0)$.
On a évidemment le diagramme commutatif d'homomorphismes canoniques
\[
  \xymatrix{
    \operatorname{Pic}(X_\lambda) \ar[r] \ar[d] & \operatorname{Pic}(X) \ar[d]\\
    \operatorname{Pic}(X'_{0\lambda}) \ar[r] & \operatorname{Pic}(X'_0)
  }
\]
Pour tout $\sh{O}_{X'_0}$-Module inversible $\sh{L}'_0$, il existe un $\lambda\geq\alpha$ et un $\sh{O}_{X'_{0\lambda}}$-Module inversible $\sh{L}'_{0\lambda}$ tel que $\sh{L}'_0$ s'en déduise par changement de base $S\to S_\lambda$ \sref{IV.8.5.2} et \sref{IV.8.5.5};
il suffit de considérer un $\sh{O}_{X_\lambda}$-Module inversible $\sh{L}_\lambda$ tel que $\operatorname{cl}(\sh{L}'_{0\lambda})$ soit l'image de $\operatorname{cl}(\sh{L}_\lambda)$ dans $\operatorname{Pic}(X'_{0\lambda})$, puis de prendre le $\sh{O}_X$-Module inversible $\sh{L}$ déduit de $\sh{L}_\lambda$ par changement de base;
en vertu de la commutativité du diagramme précédent, $\operatorname{cl}(\sh{L}'_0)$ sera l'image de $\operatorname{cl}(\sh{L})$.

Supposons donc $A$ noethérien, donc aussi $X$, et vérifions que $X$ et $X'_0$ satisfont aux conditions de \sref{IV.21.9.11}, (ii).
On a par hypothèse $\dim(X'_0)\leq1$;
d'autre part, pour vérifier la condition 2\textsuperscript{o} de \sref{IV.21.9.11}, considérons un sous-préschéma réduit $Y$ de $X$ ayant une partie localement fermée $Y$ pour ensemble sous-jacent;
le morphisme $g:Y\to S$, restriction de $f$, étant quasi-fini en chacun des points $x_i$ de $Y\cap X_0$ ($1\leq i\leq n$), on peut appliquer \sref{IV.18.5.11}, c), et on voit que $Y$ est somme des sous-préschémas ouverts $Y_i=\Spec(\sh{O}_{Y,x_i})$ ($1\leq i\leq n$) et d'un préschéma $Y''$ tel que $Y''\cap X_0=\emp$, et en outre que les injections canoniques $j_i:Y_i\to X$ sont telles que $f\circ j_i$ soit un morphisme fini.
Comme $f$ est séparé, $j_i$ est aussi un morphisme fini, donc $Y_i$ est fermé dans $X$;
on répond donc à la question en prenant $Y'=\bigcup_{1\leq i\leq n}Y_i$.

Il reste à vérifier l'hypothèse supplémentaire de \sref{IV.21.9.11}, (ii).
Or, $X'_0$ étant une courbe séparée sur le corps $\boldsymbol{k}(s_0)$, est un $\boldsymbol{k}(s_0)$-schéma quasi-projectif (chap.~V)\footnote{Le lecteur vérifiera que \sref{IV.21.9.12} n'est pas utilisé pour prouver cette propriété au chap.~V.}, donc il existe un $\sh{O}_{X'_0}$-Module ample \sref[II]{II.5.3.1}, et cela achève la démonstration.
\end{proof}

\begin{remark}[21.9.13]
\label{IV.21.9.13}
Sous les conditions de \sref{IV.21.9.12}, supposant en outre $f$ propre, le morphisme $f$ est même \emph{projectif}:
en effet, si $\sh{L}_0$ est un $\sh{O}_{X'_0}$-Module ample, il existe, en vertu de \sref{IV.21.9.12}, un $\sh{O}_X$-Module inversible $\sh{L}$ dont l'image réciproque dans $X'_0$ est isomorphe à $\sh{L}_0$, donc est ample.
Puisque tout voisinage de $s_0$ dans $S$ est nécessairement $S$ tout entier, on déduit alors de \sref{IV.9.6.4} que $\sh{L}$ est un $\sh{O}_X$-Module ample, d'où la conclusion (\sref[II]{II.5.3.1} et \sref[II]{II.5.5.3}).
\end{remark}

\subsection{Images réciproques et images directes de cycles 1-codimensionnels}
\label{subsection:IV.21.10}

Dans un chapitre ultérieur, consacré à la théorie des intersections, seront développées de façon systématique les notions d'image inverse et d'image directe de cycles.
Dans le présent numéro, nous nous contenterons de définir ces notions dans certains cas particuliers utiles, et pour les cycles $1$-codimensionnels, ces définitions étant choisies de telle sorte qu'elles soient compatibles avec les définitions correspondantes pour les diviseurs \sref{IV.21.4} et \sref{IV.21.5}, compte tenu de l'homomorphisme $\operatorname{cyc}$ défini en \sref{IV.21.6}.

\begin{env}[21.10.1]
\label{IV.21.10.1}
\oldpage[IV-4]{289}
Soient $X$, $X'$ deux préschémas localement noethériens, $f:X'\to X$ un morphisme, $T$ une partie de $X$;
supposons que l'image par $f$ de tout point maximal de $X'$ soit un point maximal de $X$, et que, pour tout $x'\in X'^{(1)}$ (i.e.\ tel que $\dim(\sh{O}_{X',x'})=1$), on ait l'une des trois conditions suivantes pour le point $x=f(x')$:
\begin{enumerate}[label=(\roman*)]
  \item $x\notin T$;
  \item $x\in X^{(1)}$ et $\sh{O}_{X',x'}$ est un $\sh{O}_{X,x}$-module plat;
  \item l'anneau $\sh{O}_{X,x}$ est factoriel et $\mathfrak{m}_{x'}\notin\operatorname{Ass}(\sh{O}_{X',x'})$.
\end{enumerate}
Sous ces conditions, nous nous proposons, pour tout cycle $1$-codimensionnel $Z$ sur $X$ de support contenu dans $T$, de définir un cycle $1$-codimensionnel $Z'=f^*(Z)$, de sorte que $Z\mapsto f^*(Z)$ soit un homomorphisme de groupes ordonnés du sous-groupe de $\mathfrak{Z}^1(X)$ formé des cycles de support contenu dans $T$, dans le groupe ordonné $\mathfrak{Z}^1(X')$.
Pour cela, soit
\[
\label{IV.21.10.1.1}
  Z=\sum_{x\in T\cap X^{(1)}}n_x\cdot\overline{\{x\}}
  \tag{21.10.1.1}
\]
où la famille des $x\in T\cap X^{(1)}$ tels que $n_x\neq0$ est localement finie.
Pour tout $x'\in X'^{(1)}$, définissons un entier $n_{x'}$ de la façon suivante, en posant $x=f(x')$:
\begin{enumerate}[label=\arabic*\textsuperscript{o}]
  \item si $x\notin T$, on prend $n_{x'}=0$;
  \item si $x\in X^{(1)}$ et si $\sh{O}_{X',x'}$ est un $\sh{O}_{X,x}$-module plat, on sait \sref{IV.6.1.1} que $\dim(\sh{O}_{X',x'}/\mathfrak{m}_x\sh{O}_{X',x'})=0$, autrement dit $\sh{O}_{X',x'}/\mathfrak{m}_x\sh{O}_{X',x'}$ est un $\sh{O}_{X',x'}$-module de longueur finie $\lambda_{x'}$, et on prend $n_{x'}=\lambda_{x'}n_x$;
  \item si $\sh{O}_{X,x}$ est factoriel et $\mathfrak{m}_{x'}\notin\operatorname{Ass}(\sh{O}_{X',x'})$, on sait \sref{IV.21.6.9} que l'homomorphisme canonique $\operatorname{cyc}:\operatorname{Div}(\sh{O}_{X,x})\to\mathfrak{Z}^1(\Spec(\sh{O}_{X,x}))$ est bijectif, et d'autre part comme $\dim(\sh{O}_{X',x'})=1$ et $\mathfrak{m}_{x'}\notin\operatorname{Ass}(\sh{O}_{X',x'})$, $\operatorname{Ass}(\sh{O}_{X',x'})$ se compose uniquement des points maximaux de $\Spec(\sh{O}_{X',x'})$, donc l'hypothèse sur $f$ implique que $f(\operatorname{Ass}(\sh{O}_{X',x'}))\subset\operatorname{Ass}(\sh{O}_{X,x})$ et il résulte de \sref{IV.21.4.5}, (ii) que l'homomorphisme $f^*:\operatorname{Div}(\sh{O}_{X,x})\to\operatorname{Div}(\sh{O}_{X',x'})$ est défini;
enfin, $x'$ étant l'unique point fermé de $\Spec(\sh{O}_{X',x'})$, $\mathfrak{Z}^1(\Spec(\sh{O}_{X',x'}))$ est canoniquement isomorphe à $\bb{Z}$.
\end{enumerate}
On a donc un homomorphisme canonique composé
\[
\label{IV.21.10.1.2}
  \mathfrak{Z}^1(\Spec(\sh{O}_{X,x}))
  \xrightarrow{\operatorname{cyc}^{-1}}\operatorname{Div}(\sh{O}_{X,x})
  \xrightarrow{f^*}\operatorname{Div}(\sh{O}_{X',x'})
  \xrightarrow{\operatorname{cyc}}\mathfrak{Z}^1(\Spec(\sh{O}_{X',x'}))
  \xrightarrow{\sim}\bb{Z}
  \tag{21.10.1.2}
\]
Si $Z_x$ est le cycle $\sum_{y\in\Spec(\sh{O}_{X,x})\cap X^{(1)}}n_y(\overline{\{y\}}\cap\Spec(\sh{O}_{X,x}))$ sur $\Spec(\sh{O}_{X,x})$, on prend alors pour $n_{x'}$ l'\emph{image} de $Z_x$ par l'homomorphisme \sref{IV.21.10.1.2}.
\end{env}

\oldpage[IV-4]{290}
\begin{env}[21.10.2]
\label{IV.21.10.2}
Nous nous proposons de montrer que:
\begin{enumerate}[label=\emph{\Alph*)}]
  \item Lorsque deux des conditions 1\textsuperscript{o}, 2\textsuperscript{o}, 3\textsuperscript{o} de \sref{IV.21.10.1} sont \emph{simultanément} satisfaites, les valeurs correspondantes de $n_{x'}$ \emph{coïncident}.
  \item L'ensemble des $x'\in X'^{(1)}$ tels que $n_{x'}\neq0$ est \emph{localement fini} dans $X'$.
\end{enumerate}

Pour démontrer A), supposons d'abord que $x\notin T$ et que $x$ vérifie l'une des conditions 2\textsuperscript{o} ou 3\textsuperscript{o} de \sref{IV.21.10.1};
alors $n_x=0$ et si l'on est dans le cas 2\textsuperscript{o}, on a $n_{x'}=0$;
si on est dans le cas 3\textsuperscript{o}, on a $Z_x=0$ puisque $\Supp(Z)\subset T$, donc encore $n_{x'}=0$.
Il reste à considérer le cas où l'on est à la fois dans le cas 2\textsuperscript{o} et dans le cas 3\textsuperscript{o};
alors, puisque $\dim(\sh{O}_{X,x})=1$, $\sh{O}_{X,x}$ est un anneau de valuation discrète;
si $t$ est une uniformisante de cet anneau, le diviseur correspondant à $Z_x$ est $\operatorname{div}(t^{n_x})$ dans $\operatorname{Div}(\sh{O}_{X,x})$, et son image dans $\operatorname{Div}(\sh{O}_{X',x'})$ est le diviseur de $t'^{n_x}$, où $t'$ est l'élément (régulier) de $\sh{O}_{X',x'}$ image de $t$.
On peut évidemment se borner au cas où $n_x=1$, et alors la définition \sref{IV.21.6.5.1} montre que l'image de $Z_x$ par \sref{IV.21.10.1.2} est le nombre $\lambda_{x'}$, ce qui achève de prouver A).

Démontrons maintenant B).
Posons $T_0=\Supp(Z)\subset T$, $T'_0=f^{-1}(T_0)$;
il suffit de prouver que la relation $n_{x'}\neq0$ implique que $x'$ appartient à l'ensemble des points maximaux de $T'_0$, ce dernier étant localement fini dans $X'$.
Il est immédiat que l'on a nécessairement $x'\in T'_0$;
si $x'$ n'était pas maximal dans l'ensemble fermé $T'_0$, il existerait
\oldpage[IV-4]{291}
une générisation $y'$ de $x'$ dans $T'_0$, distincte de $x'$, et puisque $x'\in X'^{(1)}$, $y'$ serait nécessairement un point maximal de $X'$;
par suite $y=f(y')$ serait point maximal de $X$ par hypothèse;
mais cela est absurde puisque $y\in T_0$ et que $T_0$ est purement de codimension $1$ dans $X$.
\end{env}

\begin{env}[21.10.3]
\label{IV.21.10.3}
On peut maintenant poser
\[
  f^*(Z)=\sum_{x'\in X'^{(1)}}n_{x'}\overline{\{x'\}},
\]
la somme du second membre ayant un sens en vertu de ce qu'on a prouvé dans \sref{IV.21.10.2};
on dit que le cycle $1$-codimensionnel $f^*(Z)$ est l'\emph{image réciproque} de $Z$ par $f$.
Il est immédiat que l'application $f^*$ ainsi définie est un homomorphisme de groupes ordonnés.
En outre, si $U$ est un ouvert de $X$, $V$ un ouvert de $X'$ tel que $f(V)\subset U$ et $f':V\to U$ la restriction de $f$, il résulte aussitôt des définitions que l'on a
\[
\label{IV.21.10.3.1}
  f'^*(Z|U)=f^*(Z)|V.
  \tag{21.10.3.1}
\]

Désignons par $\Gamma_T(\mathscr{Z}^1_X)$ le plus grand sous-faisceau de groupes commutatifs de $\mathscr{Z}^1_X$ de support contenu dans $T$;
il résulte de la relation \sref{IV.21.10.3.1} que les applications $\Gamma(U,\Gamma_T(\mathscr{Z}^1_X))\to\Gamma(V,\mathscr{Z}^1_{X'})$ que l'on vient de définir, définissent un homomorphisme de faisceaux de groupes commutatifs ordonnés sur $X'$
\[
  \psi^*(\Gamma_T(\mathscr{Z}^1_X))\to\mathscr{Z}^1_{X'}
\]
où $\psi$ est l'application continue sous-jacente au morphisme $f$.
\end{env}

\begin{proposition}[21.10.4]
\label{IV.21.10.4}
Supposons vérifiées les conditions de \sref{IV.21.10.1}.
Alors, pour tout diviseur $D$ sur $X$ tel que $\Supp(D)\subset T$ et que $f^*(D)$ soit défini \sref{IV.21.4.2}, on a
\[
\label{IV.21.10.4.1}
  \operatorname{cyc}(f^*(D))=f^*(\operatorname{cyc}(D)).
  \tag{21.10.4.1}
\]
\end{proposition}

\begin{proof}
La question étant locale sur $X$, on peut se borner au cas où $X=\Spec(A)$ est affine, où $D=\operatorname{div}(t)$, où $t$ est un élément régulier non inversible de $A$, et où le sous-préschéma $Y(D)$ \sref{IV.21.2.12} n'a qu'un seul point maximal $y$, de sorte que $\operatorname{cyc}(D)=n_y\cdot\overline{\{y\}}$, où $n_y$ est la longueur de $\sh{O}_{X,y}/t_y\sh{O}_{X,y}$ \sref{IV.21.6.5.1}.
On a vu dans la démonstration de \sref{IV.21.10.2}, B) que les points $x'\in X'$ tels que $\operatorname{mult}_{x'}(f^*(\operatorname{cyc}(D)))\neq0$ sont des points maximaux de $f^{-1}(\overline{\{y\}})$.
Si le point $x'$ est dans le cas 3\textsuperscript{o} de \sref{IV.21.10.1}, l'égalité des multiplicités en $x'$ des deux membres de \sref{IV.21.10.4.1} résulte de la définition de $n_{x'}$ au moyen de l'homomorphisme \sref{IV.21.10.1.2}.
Supposons au contraire que $x'$ soit dans le cas 2\textsuperscript{o} de \sref{IV.21.10.1}, et soit $x=f(x')$;
puisque $x\in X^{(1)}\cap\overline{\{y\}}$, on a nécessairement $x=y$.
Remarquons maintenant que $f^*(D)=\operatorname{div}(t')$, où $t'$ est l'image de $t$ dans $\Gamma(X',\sh{O}_{X'})$, et $Y(f^*(D))=f^{-1}(Y(D))$;
la multiplicité de $f^*(D)$ au point $x'$ est donc la longueur $n'_{x'}$ du $\sh{O}_{X',x'}$-module
\[
  \sh{O}_{X',x'}/t'_{x'}\sh{O}_{X',x'}
  =(\sh{O}_{X,y}/t_y\sh{O}_{X,y})\otimes_{\sh{O}_{X,y}}\sh{O}_{X',x'};
\]
comme il résulte de \sref{IV.4.7.1} que l'on a $n'_{x'}=\lambda_{x'}n_y$, les multiplicités au point $x'$ des deux membres de \sref{IV.21.10.4.1} sont encore égales, ce qui achève la démonstration.
\end{proof}

\begin{env}[21.10.5]
\label{IV.21.10.5}
\oldpage[IV-4]{292}
Supposons maintenant que $f:X'\to X$ soit un morphisme de préschémas localement noethériens, transformant tout point maximal de $X'$ en point maximal de $X$, et supposons en outre que pour toute partie fermée rare $T$ de $X$, $f$ vérifie les conditions de \sref{IV.21.10.1};
cela signifie encore que pour tout $x'\in X'^{(1)}$, ou bien $x=f(x')$ est point maximal de $X$, ou bien $x'$ vérifie l'une des conditions (ii), (iii) de \sref{IV.21.10.1}.
Si on tient compte de ce que tout cycle $1$-codimensionnel a un support rare dans $X$, on voit que $f^*(Z)$ est défini pour tout cycle $1$-codimensionnel $Z$ sur $X$;
en vertu de \sref{IV.21.10.3.1}, on a donc défini ainsi un homomorphisme de faisceaux de groupes commutatifs ordonnés
\[
  \psi^*(\mathscr{Z}^1_X)\to\mathscr{Z}^1_{X'}.
\]
Si de plus, pour tout diviseur $D$ sur $X$, $f^*(D)$ est défini \sref{IV.21.4.5}, le fait que le support de $D$ soit rare dans $X$ \sref{IV.21.6.6} entraîne que \sref{IV.21.10.4} est applicable, et on a donc la formule \sref{IV.21.10.4.1} pour tout diviseur $D$ sur $X$.
En particulier:
\end{env}

\begin{proposition}[21.10.6]
\label{IV.21.10.6}
Soient $X$, $X'$ deux préschémas localement noethériens, $f:X'\to X$ un morphisme plat.
Alors $f^*(Z)$ est défini pour tout cycle $1$-codimensionnel $Z$ sur $X$, $f^*(D)$ est défini pour tout diviseur $D$ sur $X$ et l'on a la relation \sref{IV.21.10.4.1}.
\end{proposition}

\begin{proof}
En effet, si $x'\in X'^{(1)}$ est tel que $x=f(x')$ ne soit pas maximal, il résulte de \sref{IV.6.1.1} que l'on a nécessairement $x\in X^{(1)}$, donc on est dans le cas (ii) de \sref{IV.21.10.1}.
On peut donc appliquer \sref{IV.21.10.5}, compte tenu de \sref{IV.21.4.5} et de \sref{IV.2.3.4}.
\end{proof}

\begin{remark}[21.10.7]
\label{IV.21.10.7}
L'existence de $f^*(Z)$ pour tout cycle $1$-codimensionnel $Z$ sur $X$ résulte déjà de l'hypothèse que $f$ est \emph{plat aux points $x'$ de $X'$ de codimension $\leq1$ dans $X'$} (i.e.\ tels que $\dim\sh{O}_{X',x'}\leq1$);
en effet, pour tout point maximal $z'\in X'$, il résulte de \sref{IV.6.1.1} que $z=f(z')$ est point maximal de $X$ puisque $\dim\sh{O}_{X,z}\leq\dim\sh{O}_{X',z'}=0$.
De même, si $x'\in X'^{(1)}$, $x=f(x')$ est maximal ou appartient à $X^{(1)}$ par \sref{IV.6.1.1};
on peut donc encore appliquer \sref{IV.21.10.5}.
\end{remark}

\begin{proposition}[21.10.8]
\label{IV.21.10.8}
Soient $X$, $X'$, $X''$ trois préschémas localement noethériens, $f:X'\to X$, $g:X''\to X'$ deux morphismes;
on suppose que $f$ (resp.\ $g$) est plat en tout point de codimension $\leq1$ dans $X'$ (resp.\ $X''$).
Alors $f\circ g$ est plat en tout point de codimension $\leq1$ dans $X''$ et pour tout cycle $1$-codimensionnel $Z$ sur $X$, on a $(f\circ g)^*(Z)=g^*(f^*(Z))$.
\end{proposition}

\begin{proof}
La première assertion résulte de \sref{IV.6.1.1} et \sref{IV.2.1.6}.
La seconde résulte de ce que $f$ (resp.\ $g$, resp.\ $f\circ g$) transforme les points maximaux de $X'$ (resp.\ $X''$, resp.\ $X''$) en points maximaux de $X$ (resp.\ $X'$, resp.\ $X$) et de ce que, si $x''\in X''^{(1)}$ est tel que $x'=g(x'')\in X'^{(1)}$ et $x=f(x')\in X^{(1)}$, on a
\[
  \operatorname{long}(\sh{O}_{X'',x''}/\mathfrak{m}_x\sh{O}_{X'',x''})
  =\operatorname{long}(\sh{O}_{X'',x''}/\mathfrak{m}_{x'}\sh{O}_{X'',x''})
   \cdot\operatorname{long}(\sh{O}_{X',x'}/\mathfrak{m}_x\sh{O}_{X',x'}).
\]
En effet, si l'on pose $A=\sh{O}_{X,x}$, $A'=\sh{O}_{X',x'}$, $A''=\sh{O}_{X'',x''}$, $\mathfrak{m}=\mathfrak{m}_x$, $\mathfrak{m}'=\mathfrak{m}_{x'}$, on a $A''/\mathfrak{m}A''=(A'/\mathfrak{m}A')\otimes_{A'}A''$, et la formule
\[
  \operatorname{long}_{A''}(A''/\mathfrak{m}A'')
  =\operatorname{long}_{A'}(A'/\mathfrak{m}A')
   \cdot\operatorname{long}_{A''}(A''/\mathfrak{m}'A'')
\]
résulte de \sref{IV.4.7.1} et de l'hypothèse de platitude de $A''$ sur $A'$.
\end{proof}

\begin{env}[21.10.9]
\label{IV.21.10.9}
Soient $X$ un préschéma localement noethérien, $\Lambda$ un groupe commutatif, noté additivement, et considéré donc comme $\bb{Z}$-module;
on notera encore par $\bb{Z}$ et $\Lambda$ les faisceaux simples sur $X$ associés aux préfaisceaux constants égaux respectivement à $\bb{Z}$ et $\Lambda$ \sref[0]{0.3.6}.
On appelle \emph{faisceau des germes de cycles $1$-codimensionnels à coefficients dans $\Lambda$} le faisceau de groupes commutatifs $\mathscr{Z}^1_X\otimes_{\bb{Z}}\Lambda$.
Si $\Lambda=\bb{Q}$, on dit que les sections
\oldpage[IV-4]{293}
de $\mathscr{Z}^1_X\otimes_{\bb{Z}}\bb{Q}$ au-dessus de $X$ sont les \emph{cycles $1$-codimensionnels à coefficients rationnels}.
Comme les fibres de $\mathscr{Z}^1_X$ sont des $\bb{Z}$-modules sans torsion \sref{IV.21.6.3}, l'homomorphisme canonique $\mathscr{Z}^1_X\to\mathscr{Z}^1_X\otimes_{\bb{Z}}\bb{Q}$ est \emph{injectif}, de sorte que les cycles $1$-codimensionnels s'identifient à des cycles $1$-codimensionnels à coefficients rationnels.
\end{env}

\begin{env}[21.10.10]
\label{IV.21.10.10}
Nous allons voir que sous certaines conditions, on peut élargir la définition de $f^*(Z)$ donnée dans \sref{IV.21.10.3} pour un cycle $1$-codimensionnel $Z$ sur $X$, mais à condition de prendre pour $f^*(Z)$ un cycle $1$-codimensionnel à coefficients rationnels sur $X'$.
Le cas plus général où nous nous plaçons est celui où $f$ transforme tout point maximal de $X'$ en un point maximal de $X$, et où, en tout point $x'\in X'^{(1)}$, on a l'une des conditions (i), (ii), (iii) de \sref{IV.21.10.1} \emph{ou une quatrième condition} (en posant $x=f(x')$):
\begin{enumerate}
  \item[(iv)] $x\in X^{(1)}$, $\mathfrak{m}_x\notin\operatorname{Ass}(\sh{O}_{X,x})$, et \emph{en outre, si l'on pose} $A=\widehat{\sh{O}}_{X,x}$, $A'=\widehat{\sh{O}}_{X',x'}$, \emph{et si $K$ désigne l'anneau total des fractions de $A$, alors $A'$ est une $A$-algèbre finie et $K'=A'\otimes_AK$ est un $K$-module} libre.
\end{enumerate}

Notons alors $r_{x'}$ le rang du $K$-module libre $K'$, $q_{x'}$ le degré de $\boldsymbol{k}(x')=\boldsymbol{k}(A')$ sur $\boldsymbol{k}(x)=\boldsymbol{k}(A)$, et posons $\mu_{x'}=r_{x'}/q_{x'}$.
Pour un cycle $1$-codimensionnel de support contenu dans $T$, donné par \sref{IV.21.10.1.1} et tout $x'\in X'^{(1)}$, on définit $c_{x'}\in\bb{Q}$, comme égal au nombre $n_{x'}\in\bb{Z}$ lorsque l'on est dans l'un des cas $1\textsuperscript{o}$, $2\textsuperscript{o}$, $3\textsuperscript{o}$ de \sref{IV.21.10.1}; mais il reste ici une quatrième possibilité:
\begin{enumerate}
  \item[$4\textsuperscript{o}$] si $x'$ vérifie la condition (iv) ci-dessus, on prend $c_{x'}=\mu_{x'}n_x\in\bb{Q}$.
\end{enumerate}
\end{env}

\begin{env}[21.10.11]
\label{IV.21.10.11}
Il faut encore prouver que lorsque la condition $4\textsuperscript{o}$ de \sref{IV.21.10.10} est vérifiée \emph{en même temps} qu'une des conditions $1\textsuperscript{o}$, $2\textsuperscript{o}$, $3\textsuperscript{o}$ de \sref{IV.21.10.1}, on a $n_{x'}=c_{x'}$.
C'est évident si $x\notin T$ puisque alors $n_x=0$.
Pour étudier les deux autres cas, notons que $\mathfrak{m}_x\sh{O}_{X',x'}$ est fermé pour la topologie $\mathfrak{m}_{x'}$-préadique, donc le complété de $\sh{O}_{X',x'}/\mathfrak{m}_x\sh{O}_{X',x'}$ pour cette dernière topologie est $A'/\mathfrak{m}_xA'$.
Si l'on est à la fois dans le cas $2\textsuperscript{o}$ et le cas $4\textsuperscript{o}$, $\sh{O}_{X',x'}/\mathfrak{m}_x\sh{O}_{X',x'}$ est discret pour la topologie $\mathfrak{m}_{x'}$-préadique, donc isomorphe à $A'/\mathfrak{m}_xA'$.
Comme $A'$ est une $A$-algèbre finie et \emph{plate} $(\mathbf{0}_{\mathrm{III}},10.2.3)$, c'est un $A$-module \emph{libre} $(\mathbf{0}_{\mathrm{III}},10.1.3)$, et le rang de $A'/\mathfrak{m}_xA'$ sur $A/\mathfrak{m}_xA=\boldsymbol{k}(x)$ est égal au rang de $A'$ sur $A$, donc aussi à celui de $K'$ sur $K$.
D'autre part ce rang est aussi égal au produit de la longueur $\lambda_{x'}$ de $\sh{O}_{X',x'}/\mathfrak{m}_x\sh{O}_{X',x'}$ (en tant que $\sh{O}_{X',x'}$-module, ou de $\boldsymbol{k}(x')$-module) par le rang $[\boldsymbol{k}(x'): \boldsymbol{k}(x)]=q_{x'}$, ce qui prouve la relation $\mu_{x'}=r_{x'}/q_{x'}=\lambda_{x'}$ dans ce cas.

Supposons enfin que l'on soit à la fois dans le cas $3\textsuperscript{o}$ et le cas $4\textsuperscript{o}$.
Alors, puisque $\dim(\sh{O}_{X,x})=1$, $\sh{O}_{X,x}$ est un anneau de valuation discrète, donc régulier.
D'autre part, $\sh{O}_{X',x'}$ est de dimension $1$, et puisque $\mathfrak{m}_{x'}\notin\operatorname{Ass}(\sh{O}_{X',x'})$, $\sh{O}_{X',x'}$ est un anneau de Cohen--Macaulay $(\mathbf{0},16.4.6)$;
enfin, puisque $A'$ est un $A$-module de type fini, $A'/\mathfrak{m}_xA'$ est un $\boldsymbol{k}(x)$-espace vectoriel de rang fini;
puisque $\sh{O}_{X',x'}/\mathfrak{m}_x\sh{O}_{X',x'}$ est contenu dans $A'/\mathfrak{m}_xA'$, c'est aussi un $\boldsymbol{k}(x)$-espace vectoriel de rang fini, donc un anneau artinien.
L'application de \sref{IV.6.1.5} montre alors que $\sh{O}_{X',x'}$ est un $\sh{O}_{X,x}$-module \emph{plat}, donc on est aussi dans le cas $2\textsuperscript{o}$, et on conclut par ce qui a été vu plus haut.

Cela étant, dans le cas que l'on a considéré, on posera
\[
  f^*(Z)=\sum_{x'\in X'^{(1)}}c_{x'}\cdot\overline{\{x'\}}.
\]
\oldpage[IV-4]{294}
qui est donc un cycle $1$-codimensionnel sur $X'$ à coefficients \emph{rationnels}.
On a encore défini ainsi un homomorphisme $f^*$ de groupes ordonnés, vérifiant \sref{IV.21.10.3.1}, et par suite un homomorphisme de faisceaux de groupes commutatifs
\[
  \psi^*(\Gamma_T(\mathscr{Z}^1_X))\to\mathscr{Z}^1_{X'}\otimes_{\bb{Z}}\bb{Q}.
\]

Lorsque $f$ vérifie les conditions précédentes pour \emph{toute partie fermée $T$ rare} dans $X$, c'est-à-dire que pour tout $x'\in X'^{(1)}$, ou bien $x=f(x')$ est point maximal de $X$, ou bien $x'$ vérifie l'une des conditions (ii), (iii) ou (iv), $f^*(Z)$ est alors défini pour tout cycle $1$-codimensionnel $Z$ sur $X$, et on a défini un homomorphisme de faisceaux de groupes commutatifs ordonnés
\[
  \psi^*(\mathscr{Z}^1_X)\to\mathscr{Z}^1_{X'}\otimes_{\bb{Z}}\bb{Q},
\]
d'où, par tensorisation, un homomorphisme de faisceaux de $\bb{Q}$-espaces vectoriels ordonnés
\[
\label{IV.21.10.11.1}
  \psi^*(\mathscr{Z}^1_X\otimes_{\bb{Z}}\bb{Q})\to\mathscr{Z}^1_{X'}\otimes_{\bb{Z}}\bb{Q}.
  \tag{21.10.11.1}
\]
\end{env}

\begin{remark}[21.10.12]
\label{IV.21.10.12}
Lorsqu'on est dans la situation de \sref{IV.21.10.10}, on peut effectivement, pour des cycles $1$-codimensionnels $Z$ sur $X$, avoir pour $f^*(Z)$ des cycles $1$-codimensionnels à coefficients \emph{non entiers}; autrement dit les nombres $\mu_{x'}$ peuvent être non entiers.
On en a un exemple en prenant l'anneau intègre complet $A$ de $(\mathbf{0}_{\mathrm{IV}},6.15.11)$, (ii) et sa clôture intégrale $A'$: le point fermé $x'$ de $\Spec(A')$ vérifie la condition (iv) de \sref{IV.21.10.10} et l'on a $\mu_{x'}=\frac12$.
\end{remark}

\begin{lemma}[21.10.13]
\label{IV.21.10.13}
Soient $A$ un anneau local noethérien de dimension $1$, $t$ un élément régulier de $A$ appartenant à l'idéal maximal $\mathfrak{m}$ (ce qui implique que $\mathfrak{m}\notin\operatorname{Ass}(A)$).
\begin{enumerate}[label=(\roman*)]
  \item Pour tout $A$-module $M$ de type fini, le noyau $N_t(M)$ et le conoyau $P_t(M)$ de l'homothétie $t_M:M\to M$ de rapport $t$ sont de longueurs finies. On pose $d_t(M)=\operatorname{long}P_t(M)-\operatorname{long}N_t(M)$.
  \item Si $0\to M'\to M\to M''\to0$ est une suite exacte de $A$-modules de type fini, on a $d_t(M)=d_t(M')+d_t(M'')$.
  \item On a $d_t(M)\geq0$ pour tout $A$-module $M$ de type fini; pour que $d_t(M)=0$, il faut et il suffit que $M$ soit de longueur finie.
  \item Si $K$ est l'anneau total des fractions de $A$ et si $M\otimes_AK$ est un $K$-module libre de rang $n$, on a $d_t(M)=n\cdot d_t(A)=n\cdot\operatorname{long}(A/tA)$.
  \item Si $M$ vérifie les hypothèses de (iv) et si de plus $t$ est $M$-régulier, on a $\operatorname{long}(M/tM)=n\cdot\operatorname{long}(A/tA)$.
\end{enumerate}
\end{lemma}

\begin{proof}
(i) $\Spec(A)$ est formé du point $\mathfrak{m}$ et des idéaux premiers minimaux $\mathfrak{p}_i$;
comme par hypothèse $t\notin\mathfrak{p}_i$ pour tout $i$ (Bourbaki, \emph{Alg. comm.}, chap.~IV, \textsection~1, n\textsuperscript{o}~1, cor.~3 de la prop.~2), l'image de $t$ dans chacun des $A_{\mathfrak{p}_i}$ est inversible, et les supports des $A$-modules de type fini $N_t(M)$ et $P_t(M)$ sont donc vides ou réduits à $\mathfrak{m}$;
on en conclut $(\mathbf{0},16.1.10)$ que ces modules sont de longueur finie.

(ii) Puisque $t$ est régulier, on a une suite exacte
\[
  0\to A\xrightarrow{t_A}A\to A/tA\to0
\]
\oldpage[IV-4]{295}
et puisque $\operatorname{Tor}_i^A(M,A)=0$ pour $i\geq1$, la suite exacte des Tor donne d'une part la suite exacte
\[
  0\to\operatorname{Tor}_1^A(M,A/tA)\to M\xrightarrow{t_M}M\to M/tM\to0
\]
et d'autre part, pour $i\geq2$,
\[
  0=\operatorname{Tor}_i^A(M,A)\to\operatorname{Tor}_i^A(M,A/tA)\to\operatorname{Tor}_{i-1}^A(M,A)=0,
\]
donc on a $N_t(M)=\operatorname{Tor}_1^A(M,A/tA)$ et $\operatorname{Tor}_i^A(M,A/tA)=0$ pour $i\geq2$;
la suite exacte des Tor donne une suite exacte
\[
  0\to\operatorname{Tor}_1^A(M',A/tA)\to\operatorname{Tor}_1^A(M,A/tA)\to\operatorname{Tor}_1^A(M'',A/tA)\to
\]
\[
  M'\otimes_A(A/tA)\to M\otimes_A(A/tA)\to M''\otimes_A(A/tA)\to0
\]
et cette suite s'écrit, d'après ce qui précède,
\[
  0\to N_t(M')\to N_t(M)\to N_t(M'')\to P_t(M')\to P_t(M)\to P_t(M'')\to0,
\]
ce qui prouve (ii).

Pour prouver (iii), notons qu'il y a une suite de composition $(M_h)$ de $M$ dont les quotients $M_h/M_{h+1}$ sont isomorphes à $A/\mathfrak{m}$ ou à un des $A/\mathfrak{p}_i$ (Bourbaki, \emph{Alg. comm.}, chap.~IV, \textsection~1, n\textsuperscript{o}~4, th.~1), et pour que $M$ soit de longueur finie, il faut et il suffit que tous ces quotients soient isomorphes à $A/\mathfrak{m}$.
Tout revient donc (en vertu de (ii)) à prouver que $d_t(A/\mathfrak{m})=0$ et $d_t(A/\mathfrak{p}_i)>0$.
Or, l'image de $t$ dans $A/\mathfrak{m}$ étant $0$, on a $N_t(A/\mathfrak{m})=A/\mathfrak{m}$ et $P_t(A/\mathfrak{m})=A/\mathfrak{m}$, d'où la première assertion;
d'autre part, l'image de $t$ dans $A/\mathfrak{p}_i$ est régulière, donc $N_t(A/\mathfrak{p}_i)=0$ et $P_t(A/\mathfrak{p}_i)=A/(tA+\mathfrak{p}_i)$ qui n'est pas réduit à $0$, d'où la seconde assertion.

(iv) Il y a une base de $M\otimes_AK$ de la forme $(x_j/s)_{1\leq j\leq n}$, où $s$ est un élément régulier de $A$ et $x_j\in M$.
Considérons l'homomorphisme $u:A^n\to M$ qui transforme l'élément $e_j$ $(1\leq j\leq n)$ de la base canonique de $A^n$ en $x_j$ et montrons que $\operatorname{Ker}(u)$ et $\operatorname{Coker}(u)$ sont de \emph{longueur finie};
en effet, pour tout $i$, l'image de $s$ dans $A_{\mathfrak{p}_i}$ est inversible, et comme $K_{\mathfrak{p}_i}=A_{\mathfrak{p}_i}$, les images des $x_j/s$ dans $M_{\mathfrak{p}_i}$ forment une base de ce $A_{\mathfrak{p}_i}$-module;
donc $u_{\mathfrak{p}_i}:A_{\mathfrak{p}_i}^n\to M_{\mathfrak{p}_i}$ est bijectif.
Comme dans (i), on en conclut que les supports de $\operatorname{Ker}(u)$ et de $\operatorname{Coker}(u)$ sont contenus dans $\{\mathfrak{m}\}$, donc que ces modules sont de longueur finie.
Cela étant, il résulte de (ii) et (iii) que l'on a $d_t(M)=d_t(A^n)=n\cdot d_t(A)=n\cdot\operatorname{long}(A/tA)$ puisque $t$ est régulier.

Enfin, il est clair que (v) se déduit aussitôt de (iv), puisque alors $N_t(M)=0$.
\end{proof}

Ce lemme permet de généraliser \sref{IV.21.10.4}:

\begin{proposition}[21.10.13]
\label{IV.21.10.13bis}
Supposons que $f$ vérifie les conditions de \sref{IV.21.10.10}.
Alors, pour tout diviseur $D$ sur $X$ tel que $\Supp(D)\subset T$ et que $f^*(D)$ soit défini, on a
\[
\label{IV.21.10.13.1}
  \operatorname{cyc}(f^*(D))=f^*(\operatorname{cyc}(D)).
  \tag{21.10.13.1}
\]
\end{proposition}

\begin{proof}
Raisonnant comme dans \sref{IV.21.10.4}, tout revient à voir (avec les mêmes notations) que si $x'$ est dans le cas $4\textsuperscript{o}$ de \sref{IV.21.10.10} et si $n_y$ est la longueur de $\sh{O}_{X,y}/t_y\sh{O}_{X,y}$, alors la
\oldpage[IV-4]{296}
longueur $n_{x'}$ de $\sh{O}_{X',x'}/t'_{x'}\sh{O}_{X',x'}$ est égale à $\mu_{x'}n_y$, $\mu_{x'}$ étant le nombre rationnel défini dans \sref{IV.21.10.10}.
Comme $\sh{O}_{X',x'}/t'_{x'}\sh{O}_{X',x'}$ est de longueur finie, il a même longueur que son complété $\mathfrak{m}_{x'}$-préadique $A'/t'_{x'}A'$, qu'on peut aussi écrire $A'/t_yA'$;
d'ailleurs, comme $t'_{x'}$ est régulier par hypothèse dans $\sh{O}_{X',x'}$, il l'est aussi dans $A'$ par platitude $(\mathbf{0}_{\mathrm{I}},6.3.4)$, et lorsque $A'$ est considéré comme $A$-module, on peut aussi dire que $t_y$ est $A'$-régulier.
Comme $A$ est de dimension $1$ et que $A'$ est un $A$-module de type fini tel que $A'\otimes_AK$ soit un $K$-module libre de rang $r_{x'}$, on peut appliquer \sref{IV.21.10.13}, (v) à $M=A'$ et à $t_y$, et la longueur de $A'/t_yA'$ en tant que $A$-module est donc $r_{x'}n_y$.
Comme $\boldsymbol{k}(x')$ est un $\boldsymbol{k}(x)$-espace vectoriel de rang $q_{x'}$, la longueur de $A'/t_yA'$ en tant que $A'$-module est donc $r_{x'}n_y/q_{x'}=\mu_{x'}n_y$, ce qui achève la démonstration.
\end{proof}

\begin{env}[21.10.14]
\label{IV.21.10.14}
Soient maintenant $X$, $X'$ deux préschémas localement noethériens, $f:X'\to X$ un morphisme ayant les \emph{deux propriétés} suivantes:
\begin{enumerate}[label=\emph{\alph*)}]
  \item $f$ est fini;
  \item l'image par $f$ de tout point maximal de $X'$ est un point maximal de $X$.
\end{enumerate}
Pour tout $x\in X^{(1)}$, les points $x'\in f^{-1}(x)$ appartiennent alors tous à $X'^{(1)}$, comme il résulte de l'hypothèse \emph{b)} et de l'inégalité $(\mathbf{0},16.3.9.1)$, la fibre $f^{-1}(x)$ étant discrète.
Soit alors
\[
  Z'=\sum_{x'\in X'^{(1)}}n_{x'}\cdot\overline{\{x'\}}
\]
un cycle $1$-codimensionnel sur $X'$.
Pour tout $x\in X^{(1)}$, définissons un entier $n_x$ par la formule
\[
  n_x=\sum_{x'\in f^{-1}(x)}n_{x'}\cdot[\boldsymbol{k}(x'): \boldsymbol{k}(x)],
\]
ce qui a un sens, les points de $f^{-1}(x)$ étant en nombre fini et $\boldsymbol{k}(x')$ étant un corps de degré fini sur $\boldsymbol{k}(x)$ \sref[I]{I.6.4.4}.
En outre, l'ensemble des $x\in X^{(1)}$ tels que $n_x\neq0$ est \emph{localement fini} dans $X$, car il est contenu dans l'image par $f$ de l'ensemble des $x'\in X'^{(1)}$ tels que $n_{x'}\neq0$, et la conclusion résulte de ce que le morphisme $f$ est quasi-compact.
On peut donc définir un cycle $1$-codimensionnel sur $X$ en posant
\[
\label{IV.21.10.14.1}
  f_*(Z')=\sum_{x\in X^{(1)}}n_x\cdot\overline{\{x\}}
  \tag{21.10.14.1}
\]
et on dit que $f_*(Z')$ est l'\emph{image directe} de $Z'$ par $f$.
Il est clair que l'application $f_*:\mathfrak{Z}^1(X')\to\mathfrak{Z}^1(X)$ ainsi définie est un homomorphisme de groupes ordonnés.
En outre, si $U$ est un ouvert de $X$, $f_U:f^{-1}(U)\to U$ la restriction de $f$, il résulte aussitôt des définitions que l'on a
\[
\label{IV.21.10.14.2}
  (f_U)_*(Z'|f^{-1}(U))=f_*(Z')|U,
  \tag{21.10.14.2}
\]
donc, en désignant par $\psi$ l'application continue sous-jacente au morphisme $f$, les applications $\Gamma(f^{-1}(U),\mathscr{Z}^1_{X'})\to\Gamma(U,\mathscr{Z}^1_X)$ que l'on vient de définir définissent un homomorphisme de faisceaux de groupes commutatifs ordonnés sur $X$
\[
  \psi_*(\mathscr{Z}^1_{X'})\to\mathscr{Z}^1_X.
\]
\end{env}

\begin{proposition}[21.10.15]
\label{IV.21.10.15}
\oldpage[IV-4]{297}
Soient $X$, $X'$, $X''$ trois préschémas localement noethériens, $f:X'\to X$, $g:X''\to X'$ deux morphismes vérifiant les conditions \emph{a)} et \emph{b)} de \sref{IV.21.10.14}.
Alors $f\circ g$ vérifie les mêmes conditions, et pour tout cycle $1$-codimensionnel $Z''$ sur $X''$, on a
\[
  (f\circ g)_*(Z'')=f_*(g_*(Z'')).
\]
\end{proposition}

\begin{proof}
Cela résulte aussitôt des définitions.
\end{proof}

\begin{proposition}[21.10.16]
\label{IV.21.10.16}
Soient $X$, $X'$, $X_1$ trois préschémas localement noethériens, $f:X'\to X$ un morphisme vérifiant les conditions \emph{a)} et \emph{b)} de \sref{IV.21.10.14}, $g:X_1\to X$ un morphisme plat.
On pose $X'_1=X'\times_XX_1$ (de sorte que $X'_1$ est localement noethérien) et on note $f_1:X'_1\to X_1$ et $g_1:X'_1\to X'$ les projections canoniques.
Alors $f_1$ vérifie les conditions \emph{a)} et \emph{b)} de \sref{IV.21.10.14}, et pour tout $1$-cycle codimensionnel $Z'$ sur $X'$, on a
\[
\label{IV.21.10.16.1}
  g^*(f_*(Z'))=(f_1)_*(g_1^*(Z')).
  \tag{21.10.16.1}
\]
\end{proposition}

\begin{proof}
Il est clair que $f_1$ est fini, et il vérifie la condition \emph{b)} de \sref{IV.21.10.14} en vertu de \sref{IV.2.3.7}.
Pour prouver \sref{IV.21.10.16.1}, on est aussitôt ramené au cas où $X$, $X'$ et $X_1$ sont des spectres d'anneaux locaux noethériens de dimension $1$, $A$, $A'$ et $A_1$, $A'$ étant un $A$-module fini et $A_1$ un $A$-module plat.
Désignant par $x$, $x'$ et $x_1$ les points fermés de $X$, $X'$ et $X_1$ respectivement, il s'agit de voir que l'on a
\[
\label{IV.21.10.16.2}
  \sum_{x'_1}\lambda_{x'_1}[\boldsymbol{k}(x'_1):\boldsymbol{k}(x_1)]
  =\lambda_{x_1}[\boldsymbol{k}(x'):\boldsymbol{k}(x)]
  \tag{21.10.16.2}
\]
où $x'_1$ parcourt l'ensemble des points fermés de $X'_1$ (i.e. l'ensemble des points à la fois au-dessus de $x'$ et de $x_1$) et où $\lambda_{x'_1}=\operatorname{long}(\sh{O}_{X'_1,x'_1}/\mathfrak{m}_{x'}\sh{O}_{X'_1,x'_1})$ et $\lambda_{x_1}=\operatorname{long}(A_1/\mathfrak{m}_xA_1)$.
Comme on a
\[
  [\boldsymbol{k}(x'_1):\boldsymbol{k}(x')]\cdot[\boldsymbol{k}(x'):\boldsymbol{k}(x)]
  =[\boldsymbol{k}(x'_1):\boldsymbol{k}(x_1)]\cdot[\boldsymbol{k}(x_1):\boldsymbol{k}(x)],
\]
le premier membre de \sref{IV.21.10.16.2} s'écrit aussi
\[
  \frac{[\boldsymbol{k}(x'): \boldsymbol{k}(x)]}{[\boldsymbol{k}(x_1): \boldsymbol{k}(x)]}
  \operatorname{long}_{A'}(A'_1/\mathfrak{m}_{x'}A'_1)
  =\operatorname{long}_{A_1}(A'_1/\mathfrak{m}_{x'}A'_1),
\]
où on a posé $A'_1=A'\otimes_AA_1$.
On a donc $A'_1/\mathfrak{m}_{x'}A'_1=(A'/\mathfrak{m}_{x'}A')\otimes_AA_1$, et comme $A_1$ est un $A$-module plat, on a par \sref{IV.4.7.1}
\[
  \operatorname{long}_{A_1}(A'_1/\mathfrak{m}_{x'}A'_1)
  =\operatorname{long}_{A'}(A'/\mathfrak{m}_{x'}A')\operatorname{long}_{A_1}(A_1/\mathfrak{m}_xA_1)
  =[\boldsymbol{k}(x'): \boldsymbol{k}(x)]\cdot\lambda_{x_1},
\]
ce qui achève la démonstration.
\end{proof}

\begin{proposition}[21.10.17]
\label{IV.21.10.17}
Soient $X$, $X'$ deux préschémas localement noethériens, $f:X'\to X$ un morphisme fini localement libre.
Alors $f$ vérifie la condition \emph{b)} de \sref{IV.21.10.14}, et pour tout diviseur $D'$ sur $X'$, on a
\[
\label{IV.21.10.17.1}
  f_*(\operatorname{cyc}(D'))=\operatorname{cyc}(f_*(D')).
  \tag{21.10.17.1}
\]
\end{proposition}

\begin{proof}
Comme $f$ est plat et fini, la condition \emph{b)} de \sref{IV.21.10.14} et la relation $f(X'^{(1)})\subset X^{(1)}$ résultent de \sref{IV.6.1.1}.
La définition \sref{IV.21.10.14.1} montre qu'on peut se ramener au cas où $X=\Spec(A)=\Spec(\sh{O}_{X,x})$ pour un $x\in X^{(1)}$;
alors $X'=\Spec(A')$, où $A'$ est une
\oldpage[IV-4]{298}
$A$-algèbre qui est un $A$-module libre de rang fini;
en outre, on peut supposer que $D'=\operatorname{div}(t')$, où $t'$ est un élément régulier de $A'$;
on a alors $f_*(D')=\operatorname{div}(t)$, où $t=\mathrm{N}_{A'/A}(t')$ est un élément régulier de $A$ \sref{IV.21.5.2}.
On peut se borner au cas où $t'$ n'est pas inversible dans $A'$, ce qui équivaut au fait que $t$ ne l'est pas dans $A$;
l'anneau $A/tA$ est alors de longueur finie et $\neq0$, et $A'/tA'$ est composé direct d'anneaux locaux artiniens $A'_i$ $(1\leq i\leq r)$, le corps résiduel de $A'_i$ étant $\boldsymbol{k}(x'_i)$, où les $x'_i$ $(1\leq i\leq r)$ sont les points de $X'$ au-dessus de $x$.
Si $t'_i$ $(1\leq i\leq r)$ est l'image de $t'$ dans $A'_i$, $A'/t'A'$ est composé direct des $A'_i/t'_iA'_i$;
comme le produit $\operatorname{long}_{A'_i}(A'_i/t'_iA'_i)\cdot[\boldsymbol{k}(x'_i):\boldsymbol{k}(x)]$ est égal à $\operatorname{long}_A(A'_i/t'_iA'_i)$, on voit que la multiplicité au point $x$ du premier membre de \sref{IV.21.10.17.1} est $\operatorname{long}_A(A'/t'A')$, si bien que la formule à démontrer se réduit à
\[
\label{IV.21.10.17.2}
  \operatorname{long}_A(A'/t'A')=\operatorname{long}_A(A/tA).
  \tag{21.10.17.2}
\]

Cette relation découle du lemme plus général suivant.
\end{proof}

\begin{lemma}[21.10.17.3]
\label{IV.21.10.17.3}
Soient $A$ un anneau local noethérien de dimension $1$, $M$ un $A$-module libre de rang fini, $u$ un endomorphisme injectif de $M$.
Alors on a
\[
\label{IV.21.10.17.4}
  \operatorname{long}_A(\operatorname{Coker}u)=\operatorname{long}_A(A/(\det u)A).
  \tag{21.10.17.4}
\]
\end{lemma}

\begin{proof}
Distinguons plusieurs cas.

I) $A$ est un anneau de valuation discrète.
Soit en effet $\pi$ une uniformisante de $A$, et remarquons que $\operatorname{long}_A(A/\pi^kA)=k$;
si $n$ est le rang de $M$ et si $\pi^{m_i}$ $(1\leq i\leq n)$ sont les facteurs invariants de $u$, $\operatorname{Coker}(u)$ est somme directe des $A$-modules $A/\pi^{m_i}A$, donc a pour longueur $m=\sum_{i=1}^nm_i$ et $\det(u)$ est produit d'un élément inversible et de $\pi^m$, d'où la conclusion dans ce cas (Bourbaki, \emph{Alg.}, chap.~VII, \textsection~4, n\textsuperscript{o}~5, cor.~1 de la prop.~4).

II) $A$ est un anneau complet intègre (de dimension $1$).
On sait alors $(\mathbf{0},19.8.8)$, (ii) qu'il y a un sous-anneau $B$ de $A$, qui est un anneau de valuation discrète, tel que $B\to A$ soit un homomorphisme local faisant de $A$ un $B$-module de type fini;
comme ce $B$-module est évidemment sans torsion, il est libre (Bourbaki, \emph{Alg. comm.}, chap.~VI, \textsection~3, n\textsuperscript{o}~6, lemme~1).
Notons $M'$ l'ensemble $M$ muni de sa structure de $B$-module (libre), par $u'$ l'endomorphisme $u$ considéré comme $B$-endomorphisme.
Il résulte de I) que l'on a
\[
\label{IV.21.10.17.5}
  \operatorname{long}_B(\operatorname{Coker}u')=\operatorname{long}_B(B/(\det u')B).
  \tag{21.10.17.5}
\]
Mais on a (Bourbaki, \emph{Alg.}, chap.~VIII, \textsection~12, n\textsuperscript{o}~2, prop.~7)
\[
  \det u'=\mathrm{N}_{A/B}(\det u),
\]
donc, en appliquant \sref{IV.21.10.17.4} à l'homothétie $x\mapsto(\det u)x$ du $B$-module libre $A$, il vient
\[
  \operatorname{long}_B(B/(\det u')B)=\operatorname{long}_B(A/(\det u)A),
\]
d'où, en portant dans \sref{IV.21.10.17.5} et divisant par $[\boldsymbol{k}(A):\boldsymbol{k}(B)]$, longueur du $B$-module $\boldsymbol{k}(A)$, on obtient \sref{IV.21.10.17.4} dans le cas envisagé.

\oldpage[IV-4]{299}
III) $A$ est un anneau complet.
Notons qu'on peut supposer en outre que $\mathfrak{m}\notin\operatorname{Ass}(A)$;
en effet, puisque $\det(u)$ est un élément régulier de $A$, si l'on avait $\mathfrak{m}\in\operatorname{Ass}(A)$, on en déduirait $\det(u)\notin\mathfrak{m}$, donc $\det(u)$ serait inversible, $u$ un automorphisme de $M$, et la formule \sref{IV.21.10.17.4} devient alors triviale, les deux membres étant nuls.

Dans ce qui suit, pour un endomorphisme $v$ d'un module $N$ sur un anneau $R$, tel que $\operatorname{Ker}v$ et $\operatorname{Coker}v$ soient de longueur finie, on posera
\[
  \chi(N,v)=\operatorname{long}_R(\operatorname{Ker}(v))-\operatorname{long}_R(\operatorname{Coker}(v)).
\]
On notera que l'hypothèse sur $v$ revient à dire que le complexe
\[
  K^0:0\to N\xrightarrow{v}N\to0
\]
a ses modules de cohomologie de longueur finie et que $\chi(N,v)=\chi(H^0(K^0))$ $(\mathbf{0}_{\mathrm{III}},11.10)$.
On en déduit que si $N'$, $N$, $N''$ sont trois $R$-modules, si l'on a un diagramme commutatif
\[
  \xymatrix{
    0 \ar[r] & N' \ar[r]^r \ar[d]^{v'} & N \ar[r]^s \ar[d]^v & N'' \ar[r] \ar[d]^{v''} & 0\\
    0 \ar[r] & N' \ar[r]_r & N \ar[r]_s & N'' \ar[r] & 0
  }
\]
dont les lignes sont exactes, et si deux des trois nombres $\chi(N,v)$, $\chi(N',v')$ et $\chi(N'',v'')$ sont définis, alors il en est de même du troisième et on a
\[
\label{IV.21.10.17.6}
  \chi(N,v)=\chi(N',v')+\chi(N'',v'').
  \tag{21.10.17.6}
\]
Cela résulte en effet aussitôt de la suite exacte de cohomologie.

Enfin, si $N$ est un $R$-module de longueur finie, on a $\chi(N,v)=0$.

Avec ces notations, on a le lemme suivant.
\end{proof}

\begin{lemma}[21.10.17.7]
\label{IV.21.10.17.7}
Soient $A$ un anneau local noethérien de dimension $1$ dont l'idéal maximal $\mathfrak{m}$ est tel que $\mathfrak{m}\notin\operatorname{Ass}(A)$;
soient $\mathfrak{p}_i$ $(1\leq i\leq n)$ les idéaux premiers minimaux de $A$.
Soient $M$ un $A$-module libre de type fini, $u$ un endomorphisme de $M$ tel que $\chi(M,u)$ soit défini;
pour chaque $i$, posons $M_i=M\otimes_A(A/\mathfrak{p}_i)$ et soit $u_i$ l'endomorphisme $u\otimes1_{A/\mathfrak{p}_i}$ de $M_i$;
alors si $\chi(M_i,u_i)$ est défini pour chaque $i$, on a
\[
\label{IV.21.10.17.8}
  \chi(M,u)=\sum_{i=1}^n\operatorname{long}(A_{\mathfrak{p}_i})\cdot\chi(M_i,u_i).
  \tag{21.10.17.8}
\]
\end{lemma}

\begin{proof}
Puisque $\mathfrak{m}\notin\operatorname{Ass}(A)$, on a une décomposition primaire réduite unique $(0)=\bigcap_{1\leq i\leq n}\mathfrak{q}_i$, où l'idéal $\mathfrak{q}_i$ est $\mathfrak{p}_i$-primaire pour $1\leq i\leq n$.
Si l'on pose $M'_i=M\otimes_A(A/\mathfrak{q}_i)$, on a donc une suite exacte
\[
  0\to M\to\bigoplus_iM'_i\to M''\to0
\]
de $A$-modules, où $M''$ est de longueur finie:
en effet, en localisant la suite exacte précédente en chacun des $\mathfrak{p}_i$, on obtient $M''_{\mathfrak{p}_i}=0$, car $(\mathfrak{q}_i)_{\mathfrak{p}_i}=0$ et $(\mathfrak{q}_j)_{\mathfrak{p}_i}=A_{\mathfrak{p}_i}$ pour $j\neq i$ (Bourbaki, \emph{Alg. comm.}, chap.~IV, \textsection~2, n\textsuperscript{o}~4, prop.~6), donc $(M'_i)_{\mathfrak{p}_i}=M_{\mathfrak{p}_i}$ et $(M'_j)_{\mathfrak{p}_i}=0$;
\oldpage[IV-4]{300}
le support de $M''$ étant donc réduit à $\mathfrak{m}$, $M''$ est de longueur finie $(\mathbf{0},16.1.10)$.
Si l'on pose $u'_i=u\otimes1_{A/\mathfrak{q}_i}$, et si $u''$ est l'endomorphisme de $M''$ déduit de $\bigoplus u'_i$ par passage aux quotients, on déduit de \sref{IV.21.10.17.6} que $\chi(M,u)+\chi(M'',u'')=\sum_i\chi(M'_i,u'_i)$, et puisque $M''$ est de longueur finie, $\chi(M'',u'')=0$.
Pour prouver \sref{IV.21.10.17.8}, on peut donc se ramener au cas où $n=1$.
On désignera alors par $\mathfrak{p}$ l'unique idéal premier minimal, qui est le nilradical de $A$;
si $A_0=A/\mathfrak{p}$, $M_0=M\otimes_AA_0$, $u_0=u\otimes1_{A_0}$, il s'agit de prouver que si $\chi(M_0,u_0)$ est défini, on a
\[
\label{IV.21.10.17.9}
  \chi(M,u)=\operatorname{long}A_{\mathfrak{p}}\cdot\chi(M_0,u_0).
  \tag{21.10.17.9}
\]

Soit $\mathfrak{n}_j$ $(0\leq j\leq r)$ la « puissance symbolique $j$-ème » de $\mathfrak{p}$, image réciproque dans $A$ de l'idéal $(\mathfrak{p}A_{\mathfrak{p}})^j$ de $A_{\mathfrak{p}}$ $(1\leq j\leq r)$, avec $\mathfrak{n}_0=A$ et $\mathfrak{n}_r=0$;
posons
\[
  M_j=\mathfrak{n}_jM/\mathfrak{n}_{j+1}M\qquad(0\leq j\leq r-1),
\]
et notons $v_j$ l'endomorphisme de $M_j$ déduit de $u$ par restriction à $\mathfrak{n}_jM$ et passage aux quotients.
Nous allons d'abord montrer que chacun des nombres $\chi(M_j,v_j)$ est défini et que l'on a
\[
\label{IV.21.10.17.10}
  \chi(M,u)=\sum_j\chi(M_j,v_j).
  \tag{21.10.17.10}
\]
La première assertion entraînera la seconde, en appliquant \sref{IV.21.10.17.6} à chacune des suites exactes
\[
  0\to\mathfrak{n}_jM/\mathfrak{n}_{j+1}M\to M/\mathfrak{n}_{j+1}M\to M/\mathfrak{n}_jM\to0.
\]

Pour prouver la première assertion, on note que si $m$ est le rang du $A$-module libre $M$, $M_j$ est isomorphe à $(\mathfrak{n}_j/\mathfrak{n}_{j+1})^m$, ou encore (puisque $\mathfrak{p}$ annule chacun des quotients $\mathfrak{n}_j/\mathfrak{n}_{j+1}$), $M_j$ est un $A_0$-module isomorphe à $M_0\otimes_{A_0}(\mathfrak{n}_j/\mathfrak{n}_{j+1})$.
Désignons par $l_j$ le rang du $A_0$-module $\mathfrak{n}_j/\mathfrak{n}_{j+1}$;
comme le corps des fractions $K_0$ de $A_0$ est le corps résiduel de $A_{\mathfrak{p}}$, $l_j$ est aussi la longueur du $A_{\mathfrak{p}}$-module $(\mathfrak{p}A_{\mathfrak{p}})^j/(\mathfrak{p}A_{\mathfrak{p}})^{j+1}$.
Il y a un système de générateurs du $A_0$-module $M_j$ qui contient une base de $M_j\otimes_{A_0}K_0$;
donc il y a un $A_0$-homomorphisme
\[
  w_j:M_0^{l_j}\to M_j
\]
dont le localisé en l'idéal $(0)$ de $A_0$ est un isomorphisme, de sorte que le support de $\operatorname{Ker}(w_j)$ et de $\operatorname{Coker}(w_j)$ est réduit à l'idéal maximal $\mathfrak{m}/\mathfrak{p}$ de $A_0$;
$\operatorname{Ker}(w_j)$ et $\operatorname{Coker}(w_j)$ sont donc des $A_0$-modules de longueur finie $(\mathbf{0},16.1.10)$.
Comme par hypothèse $\chi(M_0,u_0)$ est défini, il en est de même de $\chi(M_0^{l_j},u_0^{l_j})=l_j\chi(M_0,u_0)$ et en vertu de \sref{IV.21.10.17.6} et du fait que $\operatorname{Ker}(w_j)$ et $\operatorname{Coker}(w_j)$ sont de longueur finie, on voit que $\chi(M_j,v_j)$ est défini et égal à $l_j\chi(M_0,u_0)$;
la relation \sref{IV.21.10.17.10} donne alors
\[
  \chi(M,u)=\left(\sum_jl_j\right)\chi(M_0,u_0),
\]
et en vertu d'une remarque précédente, $\sum_jl_j$ n'est autre que la longueur de $A_{\mathfrak{p}}$, ce qui achève de prouver le lemme \sref{IV.21.10.17.7}.
\end{proof}

\oldpage[IV-4]{301}
Pour appliquer ce lemme lorsque $A$ est un anneau complet et $\mathfrak{m}\notin\operatorname{Ass}(A)$, on observe que si $u$ est injectif, il en est de même des $u_i$ (avec les notations du lemme):
en effet, $\det u$ est alors un élément régulier de $A$, donc n'appartient à aucun des $\mathfrak{p}_i$, et son image $\det u_i$ dans $A/\mathfrak{p}_i$ est par suite un élément $\neq0$ de cet anneau intègre, ce qui prouve que $u_i$ est injectif.
Comme $\operatorname{Coker}(u_i)$ est image de $\operatorname{Coker}(u)$, il est aussi de longueur finie et $\chi(M_i,u_i)$ est donc défini pour tout $i$;
on a par suite la formule \sref{IV.21.10.17.8}.
D'autre part, puisque $\det(u)$ est un élément régulier de $A$, il n'est contenu dans aucun des $\mathfrak{p}_i$;
l'idéal $(\det u)A$ est donc $\mathfrak{m}$-primaire et le quotient $A/(\det u)A$ de longueur finie.
Appliquant le même raisonnement que ci-dessus à l'homothétie injective $t:\xi\mapsto(\det u)\xi$ de $A$ et à ses images $t_i=\det u_i$ dans les $A/\mathfrak{p}_i$, il vient
\[
  \chi(A,t)=\sum_i\operatorname{long}(A_{\mathfrak{p}_i})\cdot\chi(A/\mathfrak{p}_i,t_i).
\]
Mais les anneaux $A/\mathfrak{p}_i$ sont intègres et complets, et en appliquant le résultat de II) à chacun d'eux, on obtient encore la relation \sref{IV.21.10.17.4} pour $M$ et $u$.

IV) Cas général.
Posons $A'=\widehat{A}$, $M'=M\otimes_AA'$, $u'=u\otimes1_{A'}$;
on a $\det(u')=\det(u)$, et par platitude, $\operatorname{Coker}(u')=(\operatorname{Coker}(u))_{(A')}$ et $A'/(\det u)A'=(A/(\det u)A)_{(A')}$;
comme la formule \sref{IV.21.10.17.4} est vraie pour $A'$ et $u'$ en vertu de III), elle est aussi vraie pour $A$ et $u$ en vertu de \sref{IV.4.7.1}.

Ceci termine donc la démonstration de \sref{IV.21.10.17}.

\begin{proposition}[21.10.18]
\label{IV.21.10.18}
Soient $X$, $X'$ deux préschémas localement noethériens, $f:X'\to X$ un morphisme fini, transformant tout point maximal de $X'$ en un point maximal de $X$, et vérifiant pour tout $x'\in X'$ l'une des conditions (ii), (iii), (iv) de \sref{IV.21.10.10}.
On suppose en outre qu'il existe un entier $n$ tel que, pour tout point maximal $x$ de $X$, $(f_*(\sh{O}_{X'}))_x$ soit un $\sh{O}_{X,x}$-module libre de rang $n$.
Alors, pour tout cycle $1$-codimensionnel $Z$ sur $X$, on a
\[
\label{IV.21.10.18.1}
  f_*(f^*(Z))=n\cdot Z
  \tag{21.10.18.1}
\]
(« formule de projection »).
\end{proposition}

\begin{proof}
En vertu des définitions, on est aussitôt ramené au cas où $X$ est spectre d'un anneau local noethérien $A$ de dimension $1$, de point fermé $x$, et où $Z=\overline{\{x\}}$;
posons $X'=\Spec(A')$ où $A'$ est une $A$-algèbre finie, et, pour tout idéal premier minimal $\mathfrak{p}_i$ de $A$, $A'_{\mathfrak{p}_i}$ est un $A_{\mathfrak{p}_i}$-module libre de rang $n$.
Montrons qu'on peut se borner en outre au cas où $A$ est complet.
Faisons en effet le changement de base $h:Y\to X$, où $Y=\Spec(B)$, avec $B=\widehat{A}$, et posons
\[
  Y'=X'\times_XY=\Spec(B\otimes_AA'),\qquad g=f_{\{Y\}}:Y'\to Y;
\]
le morphisme $g$ est alors fini, et comme $h$ est plat, les points maximaux de $Y$ sont au-dessus de ceux de $X$;
au-dessus de chacun des $\mathfrak{p}_i$, il n'y a qu'un nombre fini d'idéaux minimaux $\mathfrak{p}_{ij}$ de $B$, et $(B\otimes_AA')_{\mathfrak{p}_{ij}}$ est un $B_{\mathfrak{p}_{ij}}$-module libre de rang $n$.
Enfin, si $f$ vérifie l'une des conditions (ii), (iii) ou (iv) de \sref{IV.21.10.10} en chacun des points $x'$ de $f^{-1}(x)$, $g$ vérifie la condition correspondante en l'unique point $y'$ de $g^{-1}(y)$ au-dessus de $x'$ (en désignant par $y$ le point fermé de $Y$);
c'est immédiat pour les conditions (ii) et (iv);
pour la condition (iii), elle implique que $A$ est un anneau de valuation discrète, donc il en est de même de $B$, et la condition
\oldpage[IV-4]{302}
$\mathfrak{m}_{y'}\notin\operatorname{Ass}(\sh{O}_{Y',y'})$ résulte, par platitude, de la condition $\mathfrak{m}_{x'}\notin\operatorname{Ass}(\sh{O}_{X',x'})$ $(\mathbf{0},3.3.1)$.
Le morphisme $g$ vérifie donc les mêmes conditions que $f$;
si l'on prouve la formule \sref{IV.21.10.18.1} pour $g$ et $\overline{\{y\}}$, la même formule sera valable pour $f$ et $\overline{\{x\}}$, en vertu de \sref{IV.21.10.16.1}.

On peut donc supposer que $A$ est \emph{complet};
alors il en est de même de $A'$ qui est donc composé direct d'anneaux locaux complets;
on peut par suite se borner au cas où $A'$ est un anneau \emph{local}, et il reste à vérifier \sref{IV.21.10.18.1} dans chacun des cas (ii), (iii), (iv) pour le point fermé $x'$ de $X'$.
Dans le cas (ii), $A'$ étant un $A$-module plat de type fini, est un $A$-module \emph{libre} $(\mathbf{0}_{\mathrm{III}},10.1.3)$, de rang $n$ en vertu de l'hypothèse.
Or, on a par définition \sref{IV.21.10.1} et \sref{IV.21.10.3}
\[
  f^*(Z)=\lambda_{x'}\cdot\overline{\{x'\}},
\]
où $\lambda_{x'}$ est la longueur du $A'$-module $A'/\mathfrak{m}_xA'$, puis
\[
  f_*(f^*(Z))=(\lambda_{x'}[\boldsymbol{k}(x'): \boldsymbol{k}(x)])\cdot\overline{\{x\}};
\]
mais $\lambda_{x'}[\boldsymbol{k}(x'): \boldsymbol{k}(x)]$ est la longueur de $A'/\mathfrak{m}_xA'=A'\otimes_A\boldsymbol{k}(x)$ en tant que $A$-module, ou encore son rang en tant que $\boldsymbol{k}(x)$-espace vectoriel, donc est égal à $n$.

Dans le cas (iii), $A$ est un anneau de valuation discrète, donc régulier, et l'hypothèse $\mathfrak{m}_{x'}\notin\operatorname{Ass}(\sh{O}_{X',x'})$ entraîne que $A'$ est un anneau de Cohen--Macaulay $(\mathbf{0},16.4.6)$;
comme $\dim(A')=\dim(A)=1$ et que $A'/\mathfrak{m}_xA'$ est un anneau artinien, il résulte de \sref{IV.6.1.5} que $A'$ est un $A$-module plat et on est ramené au cas (ii).

Dans le cas (iv), si $K$ est l'anneau total des fractions de $A$, $K'=A'\otimes_AK$ est par hypothèse un $K$-module libre de rang $n$ et par définition \sref{IV.21.10.10}, on a
\[
  f^*(Z)=\frac{n}{[\boldsymbol{k}(x'): \boldsymbol{k}(x)]}\cdot\overline{\{x'\}},
\]
d'où $f_*(f^*(Z))=n\cdot\overline{\{x\}}$.
\end{proof}

On notera que la formule \sref{IV.21.10.18.1} est applicable en particulier lorsque le morphisme $f$ est \emph{fini} et \emph{plat} et tel que pour tout point maximal $x$ de $X$, $(f_*(\sh{O}_{X'}))_x$ soit un $\sh{O}_{X,x}$-module libre de rang $n$.

\begin{corollary}[21.10.19]
\label{IV.21.10.19}
Sous les hypothèses de \sref{IV.21.10.18}, soit $D$ un diviseur sur $X$ tel que $f^*(D)$ soit défini \sref{IV.21.4.5}; alors on a
\[
\label{IV.21.10.19.1}
  f_*(\operatorname{cyc}(f^*(D)))=n\cdot\operatorname{cyc}(D).
  \tag{21.10.19.1}
\]
\end{corollary}

\begin{proof}
Cela résulte de \sref{IV.21.10.18} et \sref{IV.21.10.13bis}.
\end{proof}
